% =====================================================================
%  เล่มคู่มือการใช้งาน (User Manual) — โครงงานคอมพิวเตอร์ วิทยาลัยการคอมพิวเตอร์ มข.
%  Compile: XeLaTeX  (Overleaf: Menu > Settings > Compiler > XeLaTeX)
%
%  วิธีใช้: กรอกข้อมูลส่วนที่ 1 แล้วเขียนเนื้อหาในโฟลเดอร์ chapters/
%  ตัวอย่างในเล่มนี้ใช้ระบบ "CS KKU Project Manager" — ลบแล้วแทนด้วยระบบของตนเอง
%  อ่าน "แนวทางการเขียนเล่มคู่มือการใช้งาน" ในส่วนต้นของ PDF และกล่องแนวทางในแต่ละบทก่อนเริ่มเขียน
% =====================================================================
\documentclass[guides]{cskkumanual}
% guides  แสดงหน้าแนวทางการเขียนและกล่อง "แนวทางการเขียน" ทุกบท
%         *** ลบคำว่า guides ออกก่อนส่งเล่ม: \documentclass{cskkumanual} ***
% ตัวเลือกอื่น เช่น \documentclass[guides,plainnumbering]{cskkumanual}
%   plainnumbering    เลขหัวข้อแบบ 1. / 1.1 / 1.1.1 (ตัวอย่างหน้า 24 ของคู่มือการจัดทำเอกสาร)
%   boldheadings      หัวข้อย่อยตัวหนา   chapterfloats  เลขภาพ/ตารางแยกตามบท
%   titleasis         ไม่แปลงชื่อภาษาอังกฤษบนปกเป็นตัวพิมพ์ใหญ่

% ---------------------------------------------------------------------
% 1) ข้อมูลเล่ม
% ---------------------------------------------------------------------
\projectCode{CS 2568/01}                 % CS/IT พ.ศ./กลุ่ม ตามที่สาขากำหนด
\titleTH{ระบบจัดการเอกสารโครงงานคอมพิวเตอร์แบบออนไลน์}
\titleEN{Online Document Management System for Computer Projects}

\firstStudentID{653450001-1}
\firstStudentPrefixTH{นาย}
\firstStudentNameTH{สมชาย ใจดี}
\firstStudentPrefixEN{Mr.}
\firstStudentNameEN{Somchai Jaidee}

\secondStudentID{653450002-9}            % เว้นว่างทั้ง 5 บรรทัดถ้าทำคนเดียว
\secondStudentPrefixTH{นางสาว}
\secondStudentNameTH{สมหญิง รักเรียน}
\secondStudentPrefixEN{Miss}
\secondStudentNameEN{Somying Rakrian}

\advisorTH{ผศ.ดร.ประสิทธิ์ สร้างสรรค์}  % คำนำหน้าติดกับชื่อ ไม่เว้นวรรค
\advisorEN{Asst. Prof. Prasit Sangsan, Ph.D.}

\courseCode{CP354 772}                   % ตรวจสอบรหัสวิชาตามหลักสูตรของรุ่น
\courseName{โครงงานคอมพิวเตอร์ 2}
\semester{2}
\academicYearTH{2568}
\submitMonthTH{มีนาคม}
\submitYearTH{2569}

\systemVersion{2.0}
\systemURL{http://localhost:5173}

\addbibresource{references.bib}

\begin{document}

% ---------------------------------------------------------------------
% 2) ปกและส่วนต้น (เลขหน้า ก ข ค ...)
% ---------------------------------------------------------------------
\makecover
\frontmatter

\begin{preface}
% =====================================================================
% คำนำ — บอกว่าคู่มือนี้เป็นของระบบอะไร ใครควรอ่าน และแบ่งเนื้อหาอย่างไร
% ไม่ควรเกิน 1 หน้า (ชื่อผู้จัดทำจะพิมพ์ท้ายหน้าให้อัตโนมัติ)
% =====================================================================
\begin{guide}
คำนำบอกว่าคู่มือเป็นของระบบอะไร รุ่นใด ผู้อ่านแต่ละกลุ่มควรอ่านบทใด และช่องทางแจ้งข้อผิดพลาด
ไม่เกิน 1 หน้า ไม่ต้องเล่าความเป็นมาของโครงงาน
\end{guide}
เอกสารคู่มือการใช้งานฉบับนี้จัดทำขึ้นเพื่ออธิบายวิธีใช้งานระบบจัดการเอกสารโครงงานคอมพิวเตอร์แบบออนไลน์
เวอร์ชัน \systemversion{} ซึ่งช่วยให้นักศึกษาและอาจารย์ที่ปรึกษาจัดการข้อเสนอโครงงาน
เขียนเอกสารด้วย \LaTeX{} และติดตามการส่งงานได้ในที่เดียว

คู่มือแบ่งเนื้อหาตามผู้อ่าน 3 กลุ่ม ผู้ใช้งานทั่วไปอ่านบทที่~\ref{ch:start} ถึงบทที่~\ref{ch:advisor}
ซึ่งอธิบายการใช้งานตามบทบาทของนักศึกษา อาจารย์ที่ปรึกษา และผู้ดูแลรายวิชา
ผู้ดูแลระบบอ่านบทที่~\ref{ch:install} ความต้องการของระบบและการติดตั้ง และบทที่~\ref{ch:admin} การดูแลระบบ
ส่วนนักพัฒนาที่รับช่วงต่ออ่านบทที่~\ref{ch:dev} คู่มือสำหรับนักพัฒนา ผู้อ่านทุกกลุ่มเริ่มจากบทที่~\ref{ch:overview}
ภาพรวมของระบบ และใช้บทที่~\ref{ch:trouble} เมื่อพบปัญหา

ผู้จัดทำหวังเป็นอย่างยิ่งว่าคู่มือฉบับนี้จะช่วยให้ผู้ใช้งานเริ่มใช้ระบบได้อย่างถูกต้องและรวดเร็ว
หากพบข้อผิดพลาดหรือมีข้อเสนอแนะ สามารถแจ้งผู้จัดทำได้ตามช่องทางในหัวข้อ~\ref{sec:contact}

\end{preface}

% =====================================================================
% ประวัติการแก้ไขเอกสาร (ไม่บังคับ) — บันทึกทุกครั้งที่ปรับคู่มือให้ตรงกับระบบรุ่นใหม่
% ลบบรรทัด % =====================================================================
% ประวัติการแก้ไขเอกสาร (ไม่บังคับ) — บันทึกทุกครั้งที่ปรับคู่มือให้ตรงกับระบบรุ่นใหม่
% ลบบรรทัด % =====================================================================
% ประวัติการแก้ไขเอกสาร (ไม่บังคับ) — บันทึกทุกครั้งที่ปรับคู่มือให้ตรงกับระบบรุ่นใหม่
% ลบบรรทัด \input{chapters/revision-history} ใน main.tex ถ้าไม่ต้องการ
% =====================================================================
\begin{revisionhistory}
\revision{1.0}{20 ม.ค. 2569}{จัดทำคู่มือครั้งแรกสำหรับระบบเวอร์ชัน 1.0}{สมชาย ใจดี}
\revision{2.0}{1 มี.ค. 2569}{เพิ่มการแก้ไขเอกสารร่วมกันแบบเรียลไทม์ ระบบความเห็น การติดตั้งด้วย Docker บทการดูแลระบบ และคู่มือสำหรับนักพัฒนา}{สมหญิง รักเรียน}
\end{revisionhistory}
 ใน main.tex ถ้าไม่ต้องการ
% =====================================================================
\begin{revisionhistory}
\revision{1.0}{20 ม.ค. 2569}{จัดทำคู่มือครั้งแรกสำหรับระบบเวอร์ชัน 1.0}{สมชาย ใจดี}
\revision{2.0}{1 มี.ค. 2569}{เพิ่มการแก้ไขเอกสารร่วมกันแบบเรียลไทม์ ระบบความเห็น การติดตั้งด้วย Docker บทการดูแลระบบ และคู่มือสำหรับนักพัฒนา}{สมหญิง รักเรียน}
\end{revisionhistory}
 ใน main.tex ถ้าไม่ต้องการ
% =====================================================================
\begin{revisionhistory}
\revision{1.0}{20 ม.ค. 2569}{จัดทำคู่มือครั้งแรกสำหรับระบบเวอร์ชัน 1.0}{สมชาย ใจดี}
\revision{2.0}{1 มี.ค. 2569}{เพิ่มการแก้ไขเอกสารร่วมกันแบบเรียลไทม์ ระบบความเห็น การติดตั้งด้วย Docker บทการดูแลระบบ และคู่มือสำหรับนักพัฒนา}{สมหญิง รักเรียน}
\end{revisionhistory}


\makecontents        % สารบัญ -> สารบัญตาราง -> สารบัญภาพ -> สารบัญสมการ (ข้ามรายการที่ว่าง)

\writingguide{chapters/writing-guide}   % แสดงเฉพาะเมื่อใช้ตัวเลือก guides

% ---------------------------------------------------------------------
% 3) เนื้อหา (เลขหน้า 1 2 3 ...)
% ---------------------------------------------------------------------
\mainmatter

% ผู้อ่านทุกกลุ่ม
\chapter{ภาพรวมของระบบ}
% =====================================================================
% บทที่ 1 ภาพรวมของระบบ (ผู้อ่าน: ทุกกลุ่ม)
% =====================================================================
\label{ch:overview}

\begin{guide}
บทนี้ทำให้ผู้อ่านรู้ภายใน 2--3 หน้าว่าระบบทำอะไร ใครใช้ และควรอ่านบทใด
\begin{itemize}
  \item อธิบายระบบ 1 ย่อหน้า พร้อมภาพขั้นตอนการทำงานหลักของระบบ
  \item ตารางผู้ใช้และสิทธิ์ (บทบาท $\times$ ฟังก์ชัน)
  \item ตารางการใช้คู่มือ บอกว่าผู้อ่านแต่ละกลุ่มควรอ่านบทใด
  \item สัญลักษณ์และรูปแบบข้อความที่ใช้ในคู่มือ
\end{itemize}
ไม่ต้องเขียนความเป็นมา วัตถุประสงค์ หรือทฤษฎีซ้ำกับเล่มรายงาน
\end{guide}

\section{เกี่ยวกับระบบ}

\systemname{} (CS KKU Project Manager) เป็นเว็บแอปพลิเคชันสำหรับจัดการเอกสารโครงงานคอมพิวเตอร์
ของวิทยาลัยการคอมพิวเตอร์ มหาวิทยาลัยขอนแก่น ตั้งแต่การยื่นข้อเสนอโครงงาน การพิจารณาของอาจารย์ที่ปรึกษา
การเขียนเอกสารฉบับสมบูรณ์ด้วย \LaTeX{} ไปจนถึงการส่งงานตามกำหนด
ผู้ใช้เข้าใช้งานผ่านเว็บเบราว์เซอร์ที่ \systemurl{} หรือที่อยู่ที่ผู้ดูแลระบบแจ้ง
ขั้นตอนการทำงานหลักของระบบแสดงดังภาพที่~\ref{fig:workflow}

% แผนภาพขั้นตอน วาดด้วย TikZ (แก้ข้อความในวงเล็บปีกกาได้โดยตรง)
\begin{figure}[H]
\centering
\begin{tikzpicture}[
    node distance=10bp and 30bp,
    box/.style={draw=cskkuStep, fill=cskkuStep!8, rounded corners=3pt, align=center,
      font=\normalfont\small, minimum height=32bp, text width=62bp, inner sep=3bp},
    adv/.style={box, draw=cskkuTip, fill=cskkuTip!8},
    arr/.style={-{Stealth[length=5bp]}, line width=0.8pt, draw=black!60},
    lbl/.style={font=\normalfont\footnotesize, text=black!75, inner sep=2bp}]
  \node[box] (a) {สมัครสมาชิก};
  \node[box, right=of a] (b) {สร้างและส่ง\\ข้อเสนอโครงงาน};
  \node[adv, right=of b] (c) {อาจารย์\\พิจารณาข้อเสนอ};
  \node[box, right=of c] (d) {ระบบสร้าง\\โครงงาน};
  \node[box, below=30bp of d] (e) {เขียนเอกสาร\\ด้วย \LaTeX};
  \node[box, left=of e] (f) {Compile\\และตรวจ PDF};
  \node[box, left=of f] (g) {ส่งงาน\\ตามกำหนด};
  \node[adv, left=of g] (h) {อาจารย์ตรวจ\\และให้ความเห็น};
  \draw[arr] (a) -- (b);
  \draw[arr] (b) -- (c);
  \draw[arr] (c) -- node[lbl, above]{อนุมัติ} (d);
  \draw[arr] (d) -- (e);
  \draw[arr] (e) -- (f);
  \draw[arr] (f) -- (g);
  \draw[arr] (g) -- (h);
  % เส้นย้อนกลับ: ขึ้นเหนือแถวบน / ลงใต้แถวล่าง เพื่อไม่ทับเส้นอื่น
  \draw[arr] (c.north) -- ++(0,12bp) -| node[lbl, pos=0.25, above]{ส่งกลับแก้ไข} (b.north);
  \draw[arr] (h.south) -- ++(0,-12bp) -| node[lbl, pos=0.25, below]{แก้ไขตามความเห็น} (e.south);
\end{tikzpicture}
\caption{ขั้นตอนการทำงานหลักของระบบ (กรอบสีเขียวคือขั้นตอนของอาจารย์ที่ปรึกษา)}
\label{fig:workflow}
\end{figure}

\section{ผู้ใช้งานและสิทธิ์การใช้งาน}

ระบบแบ่งผู้ใช้งานเป็น 3 บทบาท แต่ละบทบาทใช้งานฟังก์ชันได้ตามตารางที่~\ref{tab:roles}
เครื่องหมาย \cmark{} หมายถึงใช้งานได้ และ -- หมายถึงไม่มีสิทธิ์

\begin{table}[H]
\caption{สิทธิ์การใช้งานของผู้ใช้แต่ละบทบาท}
\label{tab:roles}
\centering
\begin{tabular}{|L{6.4cm}|C{2.2cm}|C{2.2cm}|C{2.2cm}|}
\hline
\textbf{ฟังก์ชัน} & \textbf{นักศึกษา} & \textbf{อาจารย์ที่ปรึกษา} & \textbf{ผู้ดูแลรายวิชา} \\ \hline
สร้างและส่งข้อเสนอโครงงาน      & \cmark & --     & --     \\ \hline
พิจารณาอนุมัติ/ส่งกลับข้อเสนอ  & --     & \cmark & --     \\ \hline
แก้ไขเอกสาร \LaTeX{} และ Compile & \cmark & \cmark & --     \\ \hline
แสดงความเห็นในเอกสาร          & \cmark & \cmark & --     \\ \hline
ส่งงาน (Submission)            & \cmark & --     & --     \\ \hline
กำหนดวันส่งงาน                  & --     & \cmark & \cmark \\ \hline
ดูแดชบอร์ดสรุปและส่งออก CSV    & --     & --     & \cmark \\ \hline
\end{tabular}
\end{table}

\section{การใช้คู่มือ}
\label{sec:howto-read}

คู่มือแบ่งเนื้อหาตามกลุ่มผู้อ่าน ผู้อ่านเปิดเฉพาะบทที่เกี่ยวข้องกับงานของตนได้ตามตารางที่~\ref{tab:readers}

\begin{table}[H]
\caption{บทที่ผู้อ่านแต่ละกลุ่มควรอ่าน}
\label{tab:readers}
\centering
\begin{tabular}{|L{4cm}|L{\dimexpr\linewidth-8cm-6\tabcolsep-4\arrayrulewidth\relax}|L{4cm}|}
\hline
\textbf{ผู้อ่าน} & \textbf{งานที่ทำ} & \textbf{บทที่ควรอ่าน} \\ \hline
นักศึกษา & สร้างข้อเสนอโครงงาน เขียนเอกสาร และส่งงาน
  & บทที่~\ref{ch:start}, \ref{ch:student} และ~\ref{ch:trouble} \\ \hline
อาจารย์ที่ปรึกษาและผู้ดูแลรายวิชา & พิจารณาข้อเสนอ ให้ความเห็น และกำหนดวันส่งงาน
  & บทที่~\ref{ch:start}, \ref{ch:advisor} และ~\ref{ch:trouble} \\ \hline
ผู้ดูแลระบบ & ติดตั้ง ตั้งค่า สำรองข้อมูล และปรับรุ่นระบบ
  & บทที่~\ref{ch:install}, \ref{ch:admin} และ~\ref{ch:trouble} \\ \hline
นักพัฒนา & แก้ไขและพัฒนาระบบต่อ
  & บทที่~\ref{ch:install}, \ref{ch:dev} และภาคผนวก~ข \\ \hline
\end{tabular}
\end{table}

\section{สัญลักษณ์และรูปแบบที่ใช้ในคู่มือ}

คู่มือนี้ใช้รูปแบบข้อความตามตารางที่~\ref{tab:conventions} เพื่อให้ผู้อ่านแยกองค์ประกอบบนหน้าจอออกจากเนื้อความได้ง่าย

\begin{table}[H]
\caption{สัญลักษณ์และรูปแบบข้อความที่ใช้ในคู่มือ}
\label{tab:conventions}
\centering
\begin{tabular}{|C{4.4cm}|L{8.8cm}|}
\hline
\textbf{รูปแบบ} & \textbf{ความหมาย} \\ \hline
\btn{เข้าสู่ระบบ} & ปุ่มบนหน้าจอที่ต้องคลิก \\ \hline
\field{อีเมล} & ช่องกรอกข้อมูลหรือตัวเลือกในแบบฟอร์ม \\ \hline
\menu{แทรก > ตาราง} & เมนูและเมนูย่อยที่ต้องเลือกตามลำดับ \\ \hline
\keys{Ctrl+S} & แป้นพิมพ์หรือแป้นพิมพ์ลัดที่ต้องกดพร้อมกัน \\ \hline
\calloutmark{1} & หมายเลขกำกับบนภาพหน้าจอ อธิบายไว้ใต้ภาพ \\ \hline
\cmd{npm run dev} & คำสั่งหรือข้อความที่ต้องพิมพ์ตามตัวอักษร \\ \hline
\end{tabular}
\end{table}

ข้อความสำคัญจะอยู่ในกรอบ 4 แบบ ดังนี้

\begin{note}
ข้อมูลเพิ่มเติมที่ช่วยให้เข้าใจการทำงานของระบบ
\end{note}
\begin{tip}
วิธีที่ช่วยให้ใช้งานได้สะดวกหรือรวดเร็วขึ้น
\end{tip}
\begin{warning}
สิ่งที่ต้องระวัง หากไม่ปฏิบัติตามอาจทำให้ข้อมูลผิดพลาดหรือใช้งานไม่ได้
\end{warning}
\begin{important}
ข้อกำหนดที่ต้องปฏิบัติตามเสมอ เช่น กำหนดส่งงานหรือเงื่อนไขของรายวิชา
\end{important}


% ผู้ดูแลระบบ / นักพัฒนา: ความต้องการของระบบในการติดตั้ง
\chapter{ความต้องการของระบบและการติดตั้ง}
% =====================================================================
% บทที่ 2 ความต้องการของระบบและการติดตั้ง (ผู้อ่าน: ผู้ดูแลระบบ นักพัฒนา)
% =====================================================================
\label{ch:install}

\begin{guide}
ผู้อ่านบทนี้คือผู้ที่ต้องติดตั้งระบบจากซอร์สโค้ดหรือไฟล์ที่ส่ง เช่น ผู้ดูแลระบบ กรรมการ
หรือรุ่นน้องที่รับช่วงต่อ ให้มีเนื้อหาดังนี้
\begin{itemize}
  \item ความต้องการของระบบ แยกฝั่งผู้ใช้งาน (อุปกรณ์ เบราว์เซอร์ ระบบปฏิบัติการ)
        และฝั่งเครื่องแม่ข่าย (หน่วยประมวลผล หน่วยความจำ พื้นที่ดิสก์ ระบบปฏิบัติการ)
        ระบุทั้งค่าขั้นต่ำและค่าที่แนะนำ
  \item ซอฟต์แวร์ที่ต้องติดตั้งก่อน พร้อมรุ่นและคำสั่งตรวจสอบ
  \item ขั้นตอนติดตั้งเรียงตามลำดับ ทุกคำสั่งคัดลอกไปใช้ได้ และบอกผลลัพธ์ที่ควรเห็นหลังขั้นตอนสำคัญ
  \item การตั้งค่าเป็นตาราง (ชื่อ ค่าตัวอย่าง ความหมาย) และค่าที่ต้องเปลี่ยนก่อนใช้งานจริง
  \item วิธีตรวจว่าติดตั้งสำเร็จ และบัญชีเริ่มต้นสำหรับทดสอบ
\end{itemize}
แอปบนมือถือ: ระบุรุ่น Android/iOS วิธีติดตั้งไฟล์ APK หรือจากสโตร์ และสิทธิ์ที่แอปขอ \\
ระบบฮาร์ดแวร์/IoT: ใส่รายการอุปกรณ์พร้อมรุ่น แผนผังการต่อวงจร แหล่งจ่ายไฟ และการอัปโหลดเฟิร์มแวร์ \\
งานวิจัย/โมเดล AI: ระบุรุ่น Python ไลบรารี GPU และตำแหน่งดาวน์โหลดโมเดลหรือชุดข้อมูล
\end{guide}

\section{ความต้องการของระบบ}
\label{sec:requirements}

ระบบเป็นเว็บแอปพลิเคชัน ผู้ใช้งานใช้ผ่านเว็บเบราว์เซอร์โดยไม่ต้องติดตั้งโปรแกรมเพิ่ม
ส่วนเครื่องแม่ข่ายติดตั้งส่วนหลัง ฐานข้อมูล และตัวแปลเอกสาร \LaTeX{} ซึ่งทำงานในคอนเทนเนอร์ Docker

\subsection{ความต้องการฝั่งผู้ใช้งาน}

อุปกรณ์และซอฟต์แวร์ที่ผู้ใช้งานต้องมีแสดงดังตารางที่~\ref{tab:client-req}

\begin{table}[H]
\caption{ความต้องการฝั่งผู้ใช้งาน}
\label{tab:client-req}
\centering
\begin{tabular}{|L{4cm}|L{\dimexpr\linewidth-4cm-4\tabcolsep-3\arrayrulewidth\relax}|}
\hline
\textbf{รายการ} & \textbf{ข้อกำหนด} \\ \hline
เว็บเบราว์เซอร์ & Google Chrome, Microsoft Edge หรือ Mozilla Firefox รุ่นล่าสุด \\ \hline
ความละเอียดหน้าจอ & ตั้งแต่ 1366 $\times$ 768 พิกเซลขึ้นไป (แนะนำ 1920 $\times$ 1080) \\ \hline
บัญชีผู้ใช้ & อีเมลของมหาวิทยาลัย (\cmd{@kku.ac.th}) \\ \hline
การเชื่อมต่อ & อินเทอร์เน็ตตลอดการใช้งาน (การแก้ไขร่วมกันใช้ WebSocket) \\ \hline
\end{tabular}
\end{table}

\subsection{ความต้องการของเครื่องแม่ข่าย}

เครื่องแม่ข่ายสำหรับใช้งานจริงควรมีคุณสมบัติตามตารางที่~\ref{tab:server-req}

\begin{table}[H]
\caption{ความต้องการของเครื่องแม่ข่าย}
\label{tab:server-req}
\centering
\begin{tabular}{|L{4cm}|L{\dimexpr0.5\linewidth-2cm-3\tabcolsep-2\arrayrulewidth\relax}|L{\dimexpr0.5\linewidth-2cm-3\tabcolsep-2\arrayrulewidth\relax}|}
\hline
\textbf{รายการ} & \textbf{ขั้นต่ำ} & \textbf{แนะนำ} \\ \hline
ระบบปฏิบัติการ & Ubuntu 22.04 LTS & Ubuntu 24.04 LTS \\ \hline
หน่วยประมวลผล & 2 vCPU & 4 vCPU \\ \hline
หน่วยความจำ & 4 GB & 8 GB \\ \hline
พื้นที่ดิสก์ (SSD) & 20 GB & 50 GB \\ \hline
เครือข่าย & 100 Mbps & 1 Gbps \\ \hline
\end{tabular}
\end{table}

\begin{note}
อิมเมจ TeX Live ใช้พื้นที่ประมาณ 5 GB และการ Compile แต่ละครั้งใช้หน่วยความจำได้ไม่เกิน 512 MB
จำนวนผู้ใช้ที่ Compile พร้อมกันได้จึงขึ้นกับหน่วยความจำของเครื่องแม่ข่าย
\end{note}

\subsection{ซอฟต์แวร์ที่ต้องติดตั้ง}

ติดตั้งซอฟต์แวร์ในตารางที่~\ref{tab:prereq} บนเครื่องแม่ข่ายก่อนติดตั้งระบบ
ฐานข้อมูล PostgreSQL และ TeX Live ไม่ต้องติดตั้งแยก เพราะทำงานในคอนเทนเนอร์ Docker

\begin{table}[H]
\caption{ซอฟต์แวร์ที่ต้องติดตั้งบนเครื่องแม่ข่าย}
\label{tab:prereq}
\centering
\begin{tabular}{|L{4.2cm}|L{2.6cm}|L{\dimexpr\linewidth-11cm-8\tabcolsep-5\arrayrulewidth\relax}|L{4.2cm}|}
\hline
\textbf{ซอฟต์แวร์} & \textbf{รุ่น} & \textbf{ใช้สำหรับ} & \textbf{คำสั่งตรวจสอบ} \\ \hline
Git & 2.30 ขึ้นไป & ดาวน์โหลดซอร์สโค้ด & \cmd{git --version} \\ \hline
Docker และ Docker Compose & 24 ขึ้นไป & ฐานข้อมูล ส่วนหลัง และ TeX Live & \cmd{docker --version} \\ \hline
Node.js และ npm & 20 LTS / 9 ขึ้นไป & build ส่วนหน้า & \cmd{node --version} \\ \hline
Nginx & 1.18 ขึ้นไป & เว็บเซิร์ฟเวอร์ & \cmd{nginx -v} \\ \hline
\end{tabular}
\end{table}

\section{การติดตั้งระบบ}
\label{sec:install-steps}

หัวข้อนี้ติดตั้งระบบบนเครื่องแม่ข่ายด้วย Docker Compose ซึ่งเริ่มฐานข้อมูลและส่วนหลังพร้อมกัน
ส่วนการติดตั้งบนเครื่องของนักพัฒนาอธิบายในหัวข้อ~\ref{sec:dev-setup}

\begin{steps}
  \item ดาวน์โหลดซอร์สโค้ดไว้ที่โฟลเดอร์ \filename{~/cskkuproject}
\begin{terminal}
cd ~
git clone <repo-url> cskkuproject
cd cskkuproject
\end{terminal}
  \item สร้างไฟล์ตั้งค่าของส่วนหลัง และไฟล์รหัสผ่านสำหรับ Docker Compose
        แล้วแก้ค่าตามหัวข้อ~\ref{sec:env}
\begin{terminal}
cp server/.env.example server/.env
echo "DB_PASSWORD=<รหัสผ่านฐานข้อมูล>" > .env
echo "JWT_SECRET=<คีย์ลับแบบสุ่ม>" >> .env
\end{terminal}
  \item สร้างอิมเมจ TeX Live สำหรับ Compile เอกสาร (ทำครั้งเดียว ใช้เวลานานตามความเร็วอินเทอร์เน็ต)
\begin{terminal}
docker-compose --profile build up texlive
\end{terminal}
  \item เริ่มฐานข้อมูลและส่วนหลัง แล้วตรวจสถานะ
\begin{terminal}
docker-compose up -d
docker-compose ps
\end{terminal}
        \result{บริการ \cmd{db} และ \cmd{api} มีสถานะ \cmd{Up}}
  \item สร้างข้อมูลเริ่มต้นและบัญชีทดสอบ (ครั้งแรกเท่านั้น)
\begin{terminal}
docker-compose exec api npm run db:seed
\end{terminal}
  \item Build ส่วนหน้าแล้วคัดลอกไปยังโฟลเดอร์ของเว็บเซิร์ฟเวอร์
\begin{terminal}
cd client
npm install
npm run build
sudo mkdir -p /var/www/cskkuproject
sudo cp -r dist/. /var/www/cskkuproject/
\end{terminal}
  \item ตั้งค่า Nginx ตามไฟล์ด้านล่าง แล้วโหลดการตั้งค่าใหม่ด้วยคำสั่ง \cmd{sudo nginx -s reload}
        \result{เปิดที่อยู่ของเครื่องแม่ข่ายด้วยเว็บเบราว์เซอร์แล้วพบหน้าเข้าสู่ระบบ}
\end{steps}

\begin{codefile}{/etc/nginx/sites-available/cskkuproject}
server {
    listen 80;
    server_name cskkuproject.example.com;
    root /var/www/cskkuproject;

    location / {
        try_files $uri /index.html;
    }
    location /api/ {
        proxy_pass http://127.0.0.1:3001;
    }
    location /collab {
        proxy_pass http://127.0.0.1:3001;
        proxy_http_version 1.1;
        proxy_set_header Upgrade $http_upgrade;
        proxy_set_header Connection "upgrade";
    }
}
\end{codefile}

\section{การตั้งค่าระบบ}
\label{sec:env}

ไฟล์ \filename{server/.env} กำหนดค่าการทำงานของส่วนหลังตามตารางที่~\ref{tab:env}
ส่วนไฟล์ \filename{.env} ข้างไฟล์ \filename{docker-compose.yml} เก็บรหัสผ่านฐานข้อมูล (\cmd{DB_PASSWORD})
และคีย์ลับ (\cmd{JWT_SECRET}) ที่ Docker Compose ส่งให้ส่วนหลัง

\begin{envtable}[ตัวแปรสภาพแวดล้อมในไฟล์ server/.env][tab:env]
\envvar{DATABASE_URL}{postgres://...:5432/cskkuproject}{ที่อยู่ฐานข้อมูล PostgreSQL ถ้าไม่กำหนด ระบบใช้ SQLite (สำหรับพัฒนา)}
\envvar{JWT_SECRET}{change-this-secret}{คีย์ลงนาม token สุ่มอย่างน้อย 64 ไบต์ ห้ามใช้ค่าตัวอย่างในระบบจริง}
\envvar{JWT_EXPIRES_IN}{7d}{อายุของ token หลังเข้าสู่ระบบ}
\envvar{BCRYPT_ROUNDS}{12}{จำนวนรอบการเข้ารหัสรหัสผ่าน}
\envvar{PORT}{3001}{พอร์ตของส่วนหลัง}
\envvar{NODE_ENV}{production}{\cmd{development} หรือ \cmd{production}}
\envvar{CLIENT_URL}{https://cskkuproject.example.com}{ที่อยู่ของส่วนหน้า (ใช้ตั้งค่า CORS)}
\envvar{DOCKER_IMAGE}{texlive-thai:2023}{อิมเมจที่ใช้ Compile เอกสาร}
\envvar{DOCKER_MEMORY_LIMIT}{536870912}{หน่วยความจำสูงสุดต่อการ Compile (512 MB)}
\envvar{COMPILE_TIMEOUT_MS}{60000}{เวลาสูงสุดต่อการ Compile (มิลลิวินาที)}
\envvar{MAX_FILE_SIZE}{2097152}{ขนาดไฟล์อัปโหลดสูงสุด (2 MB)}
\end{envtable}

\begin{codefile}[language=bash]{server/.env}
# Server
PORT=3001
NODE_ENV=production
CLIENT_URL=https://cskkuproject.example.com

# JWT
JWT_SECRET=change-this-to-a-strong-random-string
JWT_EXPIRES_IN=7d
BCRYPT_ROUNDS=12
\end{codefile}

\begin{tip}
สร้างค่า \cmd{JWT_SECRET} แบบสุ่มได้ด้วยคำสั่งด้านล่าง แล้วคัดลอกผลลัพธ์ไปใส่ในไฟล์ \filename{.env}
\end{tip}
\begin{terminal}
openssl rand -hex 64
\end{terminal}

\section{การตรวจสอบหลังติดตั้ง}
\label{sec:verify}

\begin{steps}
  \item ตรวจว่าส่วนหลังทำงาน
\begin{terminal}
curl http://localhost:3001/api/health
\end{terminal}
        \result{ได้ข้อความที่ขึ้นต้นด้วย \cmd{\{"status":"ok"}}
  \item เปิดระบบด้วยเว็บเบราว์เซอร์ แล้วเข้าสู่ระบบด้วยบัญชีทดสอบในตารางที่~\ref{tab:demo-accounts}
  \item เปิดโครงงานที่อนุมัติแล้ว (สร้างได้ตามหัวข้อ~\ref{sec:proposal}) แล้วคลิก \btn{Compile}
        \result{ไฟล์ PDF แสดงในพื้นที่ตัวอย่าง แสดงว่าอิมเมจ TeX Live พร้อมใช้งาน}
\end{steps}

\begin{table}[H]
\caption{บัญชีสำหรับทดสอบระบบ}
\label{tab:demo-accounts}
\centering
\begin{tabular}{|L{3.4cm}|L{\dimexpr\linewidth-6.8cm-6\tabcolsep-4\arrayrulewidth\relax}|L{3.4cm}|}
\hline
\textbf{บทบาท} & \textbf{อีเมล} & \textbf{รหัสผ่าน} \\ \hline
นักศึกษา & \cmd{somchai.stud@kku.ac.th} & \cmd{password123} \\ \hline
อาจารย์ที่ปรึกษา & \cmd{somchai@kku.ac.th} & \cmd{password123} \\ \hline
\end{tabular}
\end{table}

\begin{warning}
บัญชีในตารางที่~\ref{tab:demo-accounts} สร้างจากคำสั่ง \cmd{npm run db:seed} เพื่อการทดสอบเท่านั้น
ผู้ดูแลระบบต้องลบหรือเปลี่ยนรหัสผ่านก่อนเปิดใช้งานจริง
\end{warning}


% คู่มือผู้ใช้งานทั่วไป (หนึ่งบทต่อหนึ่งบทบาท)
\chapter{การเริ่มต้นใช้งาน}
% =====================================================================
% บทที่ 3 การเริ่มต้นใช้งาน (ผู้อ่าน: ผู้ใช้ทุกบทบาท)
%
% วิธีวางหมายเลขบนภาพ:
%   1) ใส่ตัวเลือก grid ชั่วคราว  \begin{screenshot}[grid]{images/login.png}{...}{...}
%   2) อ่านพิกัดจากเส้นตาราง (0 = ซ้าย/บน, 1 = ขวา/ล่าง) แล้วใส่ \callout{x}{y}{เลข}
%   3) ลบ grid ออกก่อนส่งเล่ม
% =====================================================================
\label{ch:start}

\begin{guide}
คู่มือผู้ใช้งานทั่วไปเริ่มที่บทนี้ อธิบายงานแรกที่ผู้ใช้ทุกบทบาทต้องทำ ได้แก่ สมัครสมาชิก เข้าสู่ระบบ
รู้จักหน้าจอหลัก และออกจากระบบ
\begin{itemize}
  \item ภาพหน้าจอหลักพร้อมหมายเลขกำกับ และคำอธิบายทุกหมายเลขใต้ภาพ
  \item เงื่อนไขที่ผู้ใช้มักพลาด เช่น รูปแบบอีเมล ความยาวรหัสผ่าน ใส่ในกล่องสำคัญ (important)
  \item ภาพที่กว้างแต่เตี้ย เช่น แถบเมนู ให้ครอปเฉพาะส่วนแล้ววางหมายเลขนอกภาพ
\end{itemize}
แอปบนมือถือ: เพิ่มการเปิดแอปครั้งแรก การอนุญาตสิทธิ์ (กล้อง ตำแหน่ง การแจ้งเตือน) และท่าทางสัมผัสที่ใช้
\end{guide}

\section{การสมัครสมาชิก}
\label{sec:register}

ผู้ใช้ใหม่ทุกบทบาทต้องสมัครสมาชิกก่อนเข้าใช้งาน โดยใช้อีเมลของมหาวิทยาลัยเท่านั้น

\begin{steps}
  \item เปิดหน้าเข้าสู่ระบบ แล้วคลิกลิงก์ \btn{สมัครสมาชิก} ใต้ปุ่มเข้าสู่ระบบ
  \item กรอกข้อมูลในแบบฟอร์มให้ครบทุกช่อง
        \begin{enumerate}
          \item \field{อีเมล} และ \field{รหัสผ่าน}
          \item \field{ชื่อ-นามสกุล (ภาษาไทย)} และ \field{ชื่อ-นามสกุล (ภาษาอังกฤษ)}
          \item \field{บทบาท} ของผู้ใช้
        \end{enumerate}
  \item ถ้าเลือกบทบาทนักศึกษา ให้กรอก \field{รหัสนักศึกษา} เช่น \cmd{663440001}
  \item คลิก \btn{สมัครสมาชิก}
        \result{ระบบสร้างบัญชีและพากลับไปยังหน้าเข้าสู่ระบบ}
\end{steps}

\begin{important}
อีเมลต้องลงท้ายด้วย \cmd{@kku.ac.th} เท่านั้น และรหัสผ่านต้องมีอย่างน้อย 6 ตัวอักษร
ระบบไม่รับอีเมลจากผู้ให้บริการอื่น เช่น Gmail
\end{important}

\section{การเข้าสู่ระบบ}
\label{sec:login}

เปิดเว็บเบราว์เซอร์ไปที่ \systemurl{} จะพบหน้าเข้าสู่ระบบดังภาพที่~\ref{fig:login}
ซึ่งประกอบด้วยส่วนต่าง ๆ ตามหมายเลขกำกับ

\begin{screenshot}[width=0.52\linewidth]{images/login.png}{หน้าเข้าสู่ระบบ}{fig:login}
  \callout{0.94}{0.49}{1}
  \callout{0.94}{0.65}{2}
  \calloutline{0.845}{0.555}{0.845}{0.625}{3}
  \highlight{0.10}{0.735}{0.80}{0.095}
  \callout{0.94}{0.78}{4}
  \callout{0.79}{0.873}{5}
\end{screenshot}

\begin{calloutlist}
  \item ช่องกรอกอีเมลของมหาวิทยาลัย
  \item ช่องกรอกรหัสผ่าน
  \item ปุ่มแสดงหรือซ่อนรหัสผ่านที่พิมพ์
  \item ปุ่มเข้าสู่ระบบ
  \item ลิงก์ไปยังหน้าสมัครสมาชิก (หัวข้อ~\ref{sec:register})
\end{calloutlist}

\begin{steps}
  \item กรอกอีเมลในช่องหมายเลข~\calloutmark{1} และรหัสผ่านในช่องหมายเลข~\calloutmark{2}
  \item คลิก \btn{เข้าสู่ระบบ} หรือกด \keys{Enter}
        \result{ระบบแสดงหน้ารายการโครงงานตามบทบาทของผู้ใช้ (หัวข้อ~\ref{sec:main-screen})}
\end{steps}

\begin{note}
หากลืมรหัสผ่าน ให้ติดต่อผู้ดูแลระบบเพื่อรีเซ็ตรหัสผ่าน (หัวข้อ~\ref{sec:contact})
\end{note}

\section{ส่วนประกอบของหน้าจอหลัก}
\label{sec:main-screen}

หลังเข้าสู่ระบบ ผู้ใช้จะเห็นหน้ารายการโครงงาน ซึ่งแบ่งเป็น 2 ส่วน ได้แก่
แถบเมนูด้านบน (ภาพที่~\ref{fig:topbar}) ที่ใช้ได้จากทุกหน้าของระบบ
และพื้นที่รายการโครงงาน (ภาพที่~\ref{fig:dashboard})

% ภาพที่กว้างแต่เตี้ย เช่น แถบเมนู: ครอปเฉพาะส่วนนั้นให้ตัวอักษรอ่านออก
% แล้ววางหมายเลขไว้ "นอกภาพ" ด้วยค่า y มากกว่า 1 พร้อมเส้นชี้ขึ้นไปยังปุ่ม
\begin{screenshot}[width=0.65\linewidth]{images/dashboard-topbar.png}{แถบเมนูด้านบน}{fig:topbar}
  \calloutline{0.088}{1.7}{0.088}{0.82}{1}
  \calloutline{0.287}{1.7}{0.287}{0.82}{2}
  \calloutline{0.527}{1.7}{0.527}{0.82}{3}
  \calloutline{0.678}{1.7}{0.678}{0.82}{4}
  \calloutline{0.843}{1.7}{0.843}{0.82}{5}
\end{screenshot}

\begin{calloutlist}
  \item การแจ้งเตือน เช่น ข้อเสนอได้รับการอนุมัติหรือมีความเห็นใหม่
  \item ชื่อผู้ใช้และบทบาทที่เข้าสู่ระบบอยู่
  \item เมนู \btn{ส่งงาน} สำหรับส่งเอกสารตามกำหนด (หัวข้อ~\ref{sec:submit})
  \item ปุ่มตั้งค่าบัญชีผู้ใช้
  \item ปุ่ม \btn{ออกจากระบบ}
\end{calloutlist}

\begin{screenshot}{images/dashboard.png}{พื้นที่รายการโครงงาน}{fig:dashboard}
  \callout{0.625}{0.155}{1}
  \callout{0.968}{0.155}{2}
  \callout{0.38}{0.314}{3}
  \calloutline{0.40}{0.51}{0.313}{0.51}{4}
  \calloutline{0.40}{0.67}{0.331}{0.67}{5}
  \calloutline{0.40}{0.83}{0.144}{0.83}{6}
\end{screenshot}

\begin{calloutlist}
  \item ช่องค้นหาโครงงานจากชื่อ
  \item ปุ่ม \btn{+ สร้างข้อเสนอ} (หัวข้อ~\ref{sec:proposal})
  \item แท็บกรองโครงงานตามสถานะ ได้แก่ ทั้งหมด ฉบับร่าง รอพิจารณา อนุมัติ และแก้ไข
  \item สถานะล่าสุดของโครงงาน
  \item การ์ดโครงงาน แสดงชื่อโครงงาน นักศึกษา และอาจารย์ที่ปรึกษา
  \item ปุ่ม \btn{ดูรายละเอียด} เปิดหน้ารายละเอียดของโครงงาน
\end{calloutlist}

\section{การออกจากระบบ}

คลิก \btn{ออกจากระบบ} ที่มุมขวาบน (หมายเลข~\calloutmark{5} ในภาพที่~\ref{fig:topbar})
ระบบบันทึกเอกสารอัตโนมัติทุก 2 วินาทีที่มีการแก้ไข แต่ควรกด \keys{Ctrl+S} ก่อนออกจากหน้าแก้ไขเอกสารทุกครั้ง


\chapter{การใช้งานสำหรับนักศึกษา}
% =====================================================================
% บทที่ 4 การใช้งานสำหรับนักศึกษา (คู่มือผู้ใช้งานทั่วไป: หนึ่งบทต่อหนึ่งบทบาท)
% =====================================================================
\label{ch:student}

\begin{guide}[แนวทางการเขียนคู่มือผู้ใช้งานทั่วไป]
เขียนหนึ่งบทต่อหนึ่งบทบาท (เช่น นักศึกษา อาจารย์ ผู้ดูแลรายวิชา) ถ้าระบบมีผู้ใช้บทบาทเดียวให้รวมเป็นบทเดียว
\begin{itemize}
  \item หนึ่งหัวข้อคือหนึ่งงาน ชื่อหัวข้อขึ้นต้นด้วย ``การ...'' เรียงตามลำดับที่ผู้ใช้ทำจริง
        ผู้อ่านจะค้นงานที่ต้องการจากสารบัญได้ทันที
  \item ทุกงานเขียนเป็นขั้นตอนด้วย \cmd{steps} และบอกผลลัพธ์ที่ควรเห็นด้วย \cmd{\result}
        ใส่ภาพหน้าจอเมื่อเป็นหน้าจอใหม่หรือหน้าจอที่มีหลายส่วน
  \item สถานะ ข้อความแจ้งเตือน และเงื่อนไข (เช่น แก้ไขได้เมื่อใด) เขียนเป็นตาราง
  \item กล่อง tip/note/warning/important ใช้เท่าที่จำเป็น ไม่เกิน 1--2 กล่องต่อหัวข้อ
  \item ไม่อธิบายการทำงานภายในของระบบ (เก็บไว้ในคู่มือสำหรับนักพัฒนา)
\end{itemize}
\end{guide}

\section{การสร้างข้อเสนอโครงงาน}
\label{sec:proposal}

นักศึกษาหัวหน้ากลุ่มเป็นผู้สร้างข้อเสนอโครงงาน แบบฟอร์มแบ่งเป็น 3 ส่วนดังภาพที่~\ref{fig:proposal}

\begin{screenshot}[width=0.5\linewidth]{images/proposal-form.png}{แบบฟอร์มรายละเอียดข้อเสนอโครงงาน}{fig:proposal}
  \callout{0.27}{0.07}{1}
  \callout{0.92}{0.237}{2}
  \callout{0.92}{0.60}{3}
  \callout{0.92}{0.76}{4}
\end{screenshot}

\begin{calloutlist}
  \item แถบสถานะของข้อเสนอ (สร้าง ส่งแล้ว รอพิจารณา)
  \item ส่วนที่ 1 ข้อมูลโครงงาน: ชื่อโครงงานภาษาไทย (จำเป็น) ชื่อภาษาอังกฤษ และรายละเอียด
  \item ส่วนที่ 2 สมาชิกโครงงาน: ผู้สร้างเป็นหัวหน้าโครงงานโดยอัตโนมัติ
  \item ส่วนที่ 3 อาจารย์ที่ปรึกษาหลัก (จำเป็น) และอาจารย์ที่ปรึกษาร่วม (ถ้ามี)
\end{calloutlist}

\begin{steps}
  \item ที่หน้ารายการโครงงาน คลิก \btn{+ สร้างข้อเสนอ}
  \item กรอก \field{ชื่อโครงงาน (ภาษาไทย)} \field{ชื่อโครงงาน (ภาษาอังกฤษ)}
        และ \field{รายละเอียด/ที่มาและความสำคัญ}
  \item เพิ่มสมาชิกโดยค้นหาจากชื่อ รหัสนักศึกษา หรืออีเมล (รวมหัวหน้าได้สูงสุด 3 คน)
  \item เลือก \field{อาจารย์ที่ปรึกษาหลัก} และ \field{อาจารย์ที่ปรึกษาร่วม} ถ้ามี
  \item คลิก \btn{บันทึกฉบับร่าง} เพื่อแก้ไขภายหลัง หรือ \btn{ส่งข้อเสนอ} เพื่อส่งให้อาจารย์พิจารณา
        \result{ข้อเสนอมีสถานะ \cmd{submitted} และอาจารย์ที่ปรึกษาได้รับการแจ้งเตือน}
\end{steps}

\begin{warning}
เลือกอาจารย์คนเดียวกันเป็นทั้งอาจารย์ที่ปรึกษาหลักและอาจารย์ที่ปรึกษาร่วมไม่ได้
และเมื่อส่งข้อเสนอแล้วจะแก้ไขไม่ได้จนกว่าอาจารย์จะส่งกลับให้แก้ไข
\end{warning}

\section{การติดตามสถานะข้อเสนอ}

ข้อเสนอโครงงานมี 4 สถานะ ตามตารางที่~\ref{tab:proposal-status}
เมื่อข้อเสนอได้รับอนุมัติ ระบบจะสร้างโครงงานพร้อมไฟล์แม่แบบ \LaTeX{} ให้โดยอัตโนมัติ

\begin{table}[H]
\caption{สถานะของข้อเสนอโครงงาน}
\label{tab:proposal-status}
\centering
\begin{tabular}{|L{2.6cm}|L{3cm}|L{7.6cm}|}
\hline
\textbf{สถานะในระบบ} & \textbf{ความหมาย} & \textbf{สิ่งที่นักศึกษาทำได้} \\ \hline
\cmd{draft}     & ฉบับร่าง       & แก้ไขและส่งข้อเสนอ \\ \hline
\cmd{submitted} & รอพิจารณา     & รอผลการพิจารณาจากอาจารย์ \\ \hline
\cmd{approved}  & อนุมัติแล้ว    & เริ่มเขียนเอกสารในโครงงานที่ระบบสร้างให้ \\ \hline
\cmd{rejected}  & ส่งกลับแก้ไข   & อ่านความเห็นของอาจารย์ แก้ไข แล้วส่งใหม่ \\ \hline
\end{tabular}
\end{table}

\section{การแก้ไขเอกสารด้วย LaTeX Editor}
\label{sec:editor}

คลิกการ์ดโครงงานที่อนุมัติแล้วเพื่อเปิด LaTeX Editor ซึ่งมีส่วนประกอบดังภาพที่~\ref{fig:editor}

\begin{screenshot}{images/editor.png}{ส่วนประกอบของหน้า LaTeX Editor}{fig:editor}
  \callout{0.53}{0.083}{1}
  \callout{0.062}{0.26}{2}
  \callout{0.34}{0.36}{3}
  \callout{0.78}{0.33}{4}
  \calloutline{0.817}{0.165}{0.817}{0.05}{5}
  \calloutline{0.863}{0.22}{0.863}{0.05}{6}
  \calloutline{0.908}{0.165}{0.908}{0.05}{7}
  \calloutline{0.963}{0.22}{0.963}{0.05}{8}
  \callout{0.22}{0.845}{9}
\end{screenshot}

\begin{calloutlist}
  \item แถบเครื่องมือ \LaTeX{} สำหรับแทรกบท หัวข้อ ตัวหนา ตาราง ภาพ และแผนภาพ Mermaid
  \item แฟ้มข้อมูลของโครงงาน สร้าง อัปโหลด และเปิดไฟล์
  \item พื้นที่แก้ไขโค้ด \LaTeX
  \item ตัวอย่าง PDF หลัง Compile
  \item ปุ่ม \btn{ซ่อน PDF} เพื่อขยายพื้นที่แก้ไข
  \item ปุ่มเปิดแผงความเห็น (หัวข้อ~\ref{sec:comment})
  \item ปุ่ม \btn{บันทึก}
  \item ปุ่ม \btn{Compile} สร้างไฟล์ PDF
  \item ผลการ Compile (COMPILE LOG) แสดงข้อผิดพลาดถ้ามี
\end{calloutlist}

แป้นพิมพ์ลัดที่ใช้บ่อยใน LaTeX Editor แสดงดังตารางที่~\ref{tab:shortcuts}

\begin{table}[H]
\caption{แป้นพิมพ์ลัดใน LaTeX Editor}
\label{tab:shortcuts}
\centering
\begin{tabular}{|L{4.6cm}|L{8.6cm}|}
\hline
\textbf{แป้นพิมพ์ลัด} & \textbf{การทำงาน} \\ \hline
\keys{Ctrl+S}       & บันทึกไฟล์ \\ \hline
\keys{Ctrl+Enter}   & Compile เอกสาร \\ \hline
\keys{Ctrl+Z}       & ย้อนกลับ (Undo) \\ \hline
\keys{Ctrl+Shift+Z} & ทำซ้ำ (Redo) \\ \hline
\keys{Ctrl+F}       & ค้นหา \\ \hline
\keys{Ctrl+H}       & ค้นหาและแทนที่ \\ \hline
\end{tabular}
\end{table}

\begin{tip}
เลือกข้อความก่อนกดปุ่ม \btn{ตัวหนา} หรือ \btn{ตัวเอียง} ระบบจะครอบข้อความด้วยคำสั่ง \LaTeX{} ให้
ถ้าไม่ได้เลือกข้อความ ระบบจะแทรกคำสั่งตัวอย่างที่ตำแหน่งเคอร์เซอร์
\end{tip}

\section{การแก้ไขเอกสารร่วมกับสมาชิก}

สมาชิกทุกคนในโครงงานเปิดไฟล์เดียวกันและแก้ไขพร้อมกันได้แบบเรียลไทม์
สถานะการเชื่อมต่อดูได้จากจุดสีที่แถบเครื่องมือ ตามตารางที่~\ref{tab:connection}

\begin{table}[H]
\caption{ความหมายของสถานะการเชื่อมต่อ}
\label{tab:connection}
\centering
% จุดสีวาดด้วย \tikz (ไม่ใช้ $\bullet$ ซึ่งระบบตรวจรูปเล่มนับเป็นสัญลักษณ์รายการ)
\begin{tabular}{|C{2.4cm}|L{10.8cm}|}
\hline
\textbf{สถานะ} & \textbf{ความหมายและสิ่งที่ควรทำ} \\ \hline
\tikz[baseline=-0.3em]\fill[cskkuTip] (0,0) circle (0.28em);~เขียว & เชื่อมต่อเรียลไทม์สำเร็จ เห็นการแก้ไขของสมาชิกคนอื่นทันที \\ \hline
\tikz[baseline=-0.3em]\fill[cskkuImp] (0,0) circle (0.28em);~แดง & ไม่ได้เชื่อมต่อ ทำงานแบบออฟไลน์และต้องบันทึกเอง \\ \hline
\tikz[baseline=-0.3em]\fill[cskkuWarn] (0,0) circle (0.28em);~เหลือง & เชื่อมต่อผิดพลาด ให้รีเฟรชหน้าเว็บ \\ \hline
\end{tabular}
\end{table}

\section{การ Compile และตรวจสอบไฟล์ PDF}

\begin{steps}
  \item คลิก \btn{Compile} หรือกด \keys{Ctrl+Enter}
  \item ตรวจผลลัพธ์ในพื้นที่ตัวอย่าง PDF (หมายเลข~\calloutmark{4} ในภาพที่~\ref{fig:editor})
  \item ถ้ามีข้อผิดพลาด อ่านรายละเอียดใน COMPILE LOG (หมายเลข~\calloutmark{9}) แล้วแก้ไขตามบรรทัดที่ระบุ
\end{steps}

\begin{note}
ระบบอนุญาตให้ Compile ได้ครั้งละ 1 คนต่อโครงงาน ถ้าเห็นข้อความ ``กำลังแปลผลอยู่''
แสดงว่าสมาชิกคนอื่นกำลัง Compile ให้รอสักครู่แล้วลองใหม่
\end{note}

\section{การส่งงาน}
\label{sec:submit}

\begin{steps}
  \item คลิกเมนู \btn{ส่งงาน} ที่แถบด้านบน
  \item คลิก \btn{สร้างการส่งงาน}
  \item เลือกโครงงานและกรอกหมายเหตุ
  \item ยืนยันการส่ง
        \result{ระบบบันทึกสำเนาของไฟล์ทั้งหมดในขณะนั้น (snapshot) พร้อมวันเวลาที่ส่ง}
\end{steps}

\begin{important}
ตรวจสอบกำหนดส่งที่หน้าแรกก่อนส่งทุกครั้ง การแก้ไขไฟล์หลังส่งจะไม่เปลี่ยนสำเนาที่ส่งไปแล้ว
\end{important}


\chapter{การใช้งานสำหรับอาจารย์ที่ปรึกษาและผู้ดูแลรายวิชา}
% =====================================================================
% บทที่ 5 การใช้งานสำหรับอาจารย์ที่ปรึกษาและผู้ดูแลรายวิชา
% =====================================================================
\label{ch:advisor}

\begin{guide}
บทของบทบาทถัดไปใช้รูปแบบเดียวกับบทที่ 4 งานที่ทำได้เหมือนกันทุกบทบาท
(เช่น การให้ความเห็น) ให้อธิบายไว้ที่เดียวแล้วอ้างถึงด้วยหมายเลขหัวข้อ
ไม่ต้องเขียนซ้ำในทุกบท
\end{guide}

\section{การพิจารณาข้อเสนอโครงงาน}

\begin{steps}
  \item ที่หน้ารายการโครงงาน เลือกแท็บ \btn{รอพิจารณา}
  \item คลิกข้อเสนอที่ต้องการเพื่อเปิดรายละเอียด
  \item พิมพ์ความเห็นในช่อง \field{ความเห็น}
  \item คลิก \btn{อนุมัติ} หรือ \btn{ส่งกลับแก้ไข}
        \result{ถ้าอนุมัติ ระบบสร้างโครงงานพร้อมไฟล์แม่แบบ 12 ไฟล์ให้นักศึกษา
        ถ้าส่งกลับ นักศึกษาจะได้รับการแจ้งเตือนให้แก้ไข}
\end{steps}

\begin{tip}
เขียนความเห็นให้ระบุสิ่งที่ต้องแก้เป็นข้อ ๆ นักศึกษาจะเห็นความเห็นนี้ที่หน้าข้อเสนอเมื่อถูกส่งกลับ
\end{tip}

\section{การให้ความเห็นในเอกสาร}
\label{sec:comment}

อาจารย์ที่ปรึกษาให้ความเห็นในไฟล์เอกสารของนักศึกษาได้โดยตรงผ่านแผงความเห็นดังภาพที่~\ref{fig:comment}
นักศึกษาตอบกลับได้ในแผงเดียวกัน

\begin{screenshot}{images/editor-comment.png}{แผงความเห็นใน LaTeX Editor}{fig:comment}
  \calloutline{0.70}{0.13}{0.60}{0.13}{1}
  \callout{0.665}{0.33}{2}
  \callout{0.665}{0.41}{3}
  \callout{0.665}{0.55}{4}
  \calloutline{0.862}{0.17}{0.862}{0.05}{5}
\end{screenshot}

\begin{calloutlist}
  \item ส่วนหัวแผงความเห็น แสดงจำนวนความเห็นและไฟล์ที่เปิดอยู่ กรองได้ตามสถานะ ทั้งหมด เปิด และแก้ไขแล้ว
  \item ช่องพิมพ์ความเห็นใหม่ เลือกข้อความใน Editor ก่อนเพื่อผูกความเห็นกับตำแหน่งนั้น
  \item ปุ่ม \btn{เพิ่มความเห็น} (หรือกด \keys{Ctrl+Enter})
  \item ความเห็นที่มีอยู่ พร้อมปุ่ม \btn{ตอบกลับ} และ \btn{แก้ไขแล้ว}
  \item ปุ่มเปิดหรือปิดแผงความเห็น ตัวเลขคือจำนวนความเห็นที่ยังเปิดอยู่
\end{calloutlist}

\begin{steps}
  \item เปิดโครงงานของนักศึกษา แล้วคลิกปุ่มหมายเลข~\calloutmark{5} เพื่อเปิดแผงความเห็น
  \item เลือกข้อความใน Editor ที่ต้องการให้ความเห็น
  \item พิมพ์ความเห็นในช่องหมายเลข~\calloutmark{2} แล้วคลิก \btn{เพิ่มความเห็น}
  \item เมื่อนักศึกษาแก้ไขตามความเห็นแล้ว คลิก \btn{แก้ไขแล้ว} เพื่อปิดความเห็น
\end{steps}

\section{การกำหนดวันส่งงาน}

อาจารย์ที่ปรึกษาและผู้ดูแลรายวิชากำหนดวันส่งงานของแต่ละรายการได้
นักศึกษาจะเห็นกำหนดส่งที่หน้าแรก และระบบจะแจ้งเตือนเมื่อใกล้ถึงกำหนด

% ตัวอย่าง: เขียนขั้นตอนตามหน้าจอจริงของระบบ พร้อมภาพหน้าจอ
% \begin{steps}
%   \item ...
% \end{steps}

\section{แดชบอร์ดสำหรับผู้ดูแลรายวิชา}
\label{sec:dashboard}

ผู้ดูแลรายวิชาและอาจารย์ที่ปรึกษาเปิดเมนูแดชบอร์ดได้ที่ \cmd{/dashboard} เพื่อดูจำนวนโครงงาน
ข้อเสนอที่รอพิจารณา และจำนวนการส่งงาน และส่งออกข้อมูลเป็นไฟล์ CSV ได้
อาจารย์ที่ต้องเป็นผู้ดูแลรายวิชาให้แจ้งผู้ดูแลระบบกำหนดบทบาทให้ (หัวข้อ~\ref{sec:users})


% คู่มือผู้ดูแลระบบ
\chapter{การดูแลระบบ}
% =====================================================================
% บทที่ 6 การดูแลระบบ (ผู้อ่าน: ผู้ดูแลระบบ)
% =====================================================================
\label{ch:admin}

\begin{guide}
บทนี้สำหรับผู้ที่ดูแลระบบหลังติดตั้ง เช่น เจ้าหน้าที่ไอทีของหน่วยงานที่รับระบบไปใช้
\begin{itemize}
  \item ตารางงานประจำและความถี่ที่ควรทำ
  \item การจัดการผู้ใช้และสิทธิ์ ถ้าระบบมีหน้าจอผู้ดูแล ให้เขียนเป็นขั้นตอนพร้อมภาพแบบคู่มือผู้ใช้
  \item การสำรองและกู้คืนข้อมูล ใส่เฉพาะคำสั่งที่ทดลองกู้คืนได้จริงแล้ว
  \item การปรับรุ่นระบบ การตรวจสถานะ และตำแหน่งบันทึกการทำงาน (log)
  \item ข้อควรระวังด้านความปลอดภัยก่อนและหลังเปิดใช้งานจริง
\end{itemize}
ระบบที่ไม่มีเครื่องแม่ข่ายหรือฐานข้อมูลกลาง เช่น โปรแกรมที่ทำงานบนเครื่องผู้ใช้อย่างเดียว ลบบทนี้ได้
แล้วย้ายเรื่องการปรับรุ่นไปไว้ในบทที่ 2
\end{guide}

\section{งานประจำของผู้ดูแลระบบ}
\label{sec:admin-tasks}

งานที่ผู้ดูแลระบบต้องทำเป็นประจำแสดงดังตารางที่~\ref{tab:admin-tasks}

\begin{table}[H]
\caption{งานประจำของผู้ดูแลระบบ}
\label{tab:admin-tasks}
\centering
\begin{tabular}{|L{\dimexpr\linewidth-7cm-6\tabcolsep-4\arrayrulewidth\relax}|L{3.6cm}|C{3.4cm}|}
\hline
\textbf{งาน} & \textbf{ความถี่} & \textbf{หัวข้อ} \\ \hline
สำรองฐานข้อมูลและไฟล์ของโครงงาน & ทุกวัน & \ref{sec:backup} \\ \hline
ตรวจสถานะระบบและพื้นที่ดิสก์ & ทุกสัปดาห์ & \ref{sec:monitor} \\ \hline
ตรวจบันทึกการทำงานเมื่อผู้ใช้แจ้งปัญหา & เมื่อมีปัญหา & \ref{sec:monitor} \\ \hline
กำหนดบทบาทผู้ดูแลรายวิชาและตั้งรหัสผ่านใหม่ให้ผู้ใช้ & เมื่อมีคำขอ & \ref{sec:users} \\ \hline
ปรับรุ่นระบบ & เมื่อมีรุ่นใหม่ & \ref{sec:upgrade} \\ \hline
\end{tabular}
\end{table}

คำสั่งในบทนี้รันที่โฟลเดอร์ของระบบบนเครื่องแม่ข่าย (\filename{~/cskkuproject}) ตามที่ติดตั้งในบทที่~\ref{ch:install}

\section{การจัดการผู้ใช้}
\label{sec:users}

ผู้ใช้สมัครสมาชิกเองด้วยอีเมลของมหาวิทยาลัย (หัวข้อ~\ref{sec:register})
ระบบมี 3 บทบาท ได้แก่ \cmd{student} (นักศึกษา) \cmd{advisor} (อาจารย์ที่ปรึกษา)
และ \cmd{instructor} (ผู้ดูแลรายวิชา) ระบบรุ่นนี้ยังไม่มีหน้าจอสำหรับผู้ดูแลระบบ
จึงจัดการผู้ใช้ด้วยคำสั่งในหัวข้อย่อยต่อไปนี้

\subsection{การดูรายชื่อผู้ใช้}

\begin{terminal}
docker-compose exec api node scripts/listUsers.js
\end{terminal}

ผลลัพธ์แสดงรหัส อีเมล บทบาท และชื่อของผู้ใช้ทุกคน หนึ่งบรรทัดต่อหนึ่งคน

\subsection{การกำหนดบทบาทผู้ดูแลรายวิชา}

อาจารย์ที่สมัครสมาชิกจะได้บทบาทตามที่เลือกตอนสมัคร
ถ้าต้องการให้อาจารย์เป็นผู้ดูแลรายวิชา ให้เปลี่ยนบทบาทเป็น \cmd{instructor} ด้วยคำสั่งด้านล่าง
(เปลี่ยนอีเมลเป็นของอาจารย์ท่านนั้น)

\begin{terminal}
docker-compose exec db psql -U cskkuadmin cskkuproject
\end{terminal}

จากนั้นพิมพ์คำสั่ง SQL ด้านล่างที่ตัวพร้อมรับคำสั่ง \cmd{cskkuproject=\#} แล้วพิมพ์ \cmd{\q} เพื่อออก

\begin{terminal}
UPDATE users SET role = 'instructor'
  WHERE email = 'teacher@kku.ac.th';
\end{terminal}

\begin{note}
บทบาทใหม่มีผลเมื่อผู้ใช้ออกจากระบบแล้วเข้าสู่ระบบอีกครั้ง เพราะบทบาทถูกบันทึกไว้ใน token ตอนเข้าสู่ระบบ
\end{note}

\subsection{การตั้งรหัสผ่านใหม่ให้ผู้ใช้}

เมื่อผู้ใช้ลืมรหัสผ่าน ผู้ดูแลระบบตั้งรหัสผ่านชั่วคราวให้ได้ด้วยคำสั่งด้านล่าง
ระบบเข้ารหัสรหัสผ่านใหม่ให้อัตโนมัติก่อนบันทึก แล้วแจ้งให้ผู้ใช้เปลี่ยนรหัสผ่านเองที่หน้าโปรไฟล์

\begin{terminal}
docker-compose exec api node -e "
const { User } = require('./models');
User.findOne({ where: { email: 'student@kku.ac.th' } })
  .then((u) => { u.password_hash = 'TempPass2569'; return u.save(); })
  .then(() => process.exit(0));"
\end{terminal}

\section{การสำรองและกู้คืนข้อมูล}
\label{sec:backup}

ข้อมูลของระบบเก็บไว้ 2 แห่ง ได้แก่ ฐานข้อมูล PostgreSQL ซึ่งเก็บผู้ใช้ ข้อเสนอ เนื้อหาไฟล์เอกสาร
ความเห็น และการส่งงาน และ volume \cmd{project_data} ซึ่งเก็บไฟล์ของโครงงานที่ใช้ Compile
ต้องสำรองทั้งสองแห่งทุกครั้ง

\subsection{การสำรองข้อมูล}

สร้างไฟล์สคริปต์ \filename{~/backup.sh} ตามด้านล่าง แล้วให้สิทธิ์รันด้วยคำสั่ง \cmd{chmod +x ~/backup.sh}
ทุกครั้งที่รันสคริปต์ จะได้ไฟล์สำรองที่มีวันที่อยู่ในชื่อไฟล์ในโฟลเดอร์ \filename{~/backup}

\begin{codefile}[language=bash]{~/backup.sh}
#!/bin/sh
mkdir -p ~/backup
cd ~/cskkuproject
# ฐานข้อมูล
docker-compose exec -T db \
  pg_dump -U cskkuadmin -Fc cskkuproject \
  > ~/backup/cskku-$(date +%F).dump
# ไฟล์ของโครงงาน
docker run --rm -v cskkuproject_project_data:/data \
  -v ~/backup:/backup alpine \
  tar czf /backup/files-$(date +%F).tar.gz -C /data .
\end{codefile}

\begin{tip}
ตั้งเวลาสำรองข้อมูลอัตโนมัติทุกคืนเวลา 02:00 น. ด้วยคำสั่ง \cmd{crontab -e} แล้วเพิ่มบรรทัด
\cmd{0 2 * * * ~/backup.sh}
\end{tip}

\subsection{การกู้คืนข้อมูล}

\begin{steps}
  \item หยุดส่วนหลังเพื่อไม่ให้มีผู้ใช้แก้ไขข้อมูลระหว่างกู้คืน
\begin{terminal}
docker-compose stop api
\end{terminal}
  \item กู้คืนฐานข้อมูลจากไฟล์สำรอง (เปลี่ยนชื่อไฟล์ตามวันที่ที่ต้องการ)
\begin{terminal}
docker-compose exec -T db \
  pg_restore -U cskkuadmin -d cskkuproject --clean \
  < ~/backup/cskku-2026-03-01.dump
\end{terminal}
  \item เริ่มส่วนหลังอีกครั้ง แล้วตรวจสอบตามหัวข้อ~\ref{sec:verify}
\begin{terminal}
docker-compose start api
\end{terminal}
\end{steps}

\begin{warning}
การกู้คืนเขียนทับข้อมูลปัจจุบันทั้งหมด ทดลองกู้คืนบนเครื่องทดสอบก่อนเสมอ
และสุ่มทดลองกู้คืนไฟล์สำรองเป็นระยะเพื่อยืนยันว่าไฟล์ใช้งานได้
\end{warning}

\section{การปรับรุ่นระบบ}
\label{sec:upgrade}

\begin{steps}
  \item สำรองข้อมูลตามหัวข้อ~\ref{sec:backup}
  \item ดาวน์โหลดซอร์สโค้ดรุ่นใหม่
\begin{terminal}
git pull
\end{terminal}
  \item สร้างและเริ่มส่วนหลังรุ่นใหม่
\begin{terminal}
docker-compose up -d --build api
\end{terminal}
  \item Build ส่วนหน้าใหม่และคัดลอกไปยังโฟลเดอร์ของเว็บเซิร์ฟเวอร์ ตามขั้นตอนที่ 6 ในหัวข้อ~\ref{sec:install-steps}
  \item ตรวจสอบระบบตามหัวข้อ~\ref{sec:verify}
\end{steps}

\begin{important}
ส่วนหลังสร้างตารางที่ยังไม่มีในฐานข้อมูลให้อัตโนมัติเมื่อเริ่มทำงาน แต่ไม่แก้ไขตารางเดิม
ถ้ารุ่นใหม่เปลี่ยนโครงสร้างตาราง ให้ทำตามบันทึกการเปลี่ยนแปลงของรุ่นนั้นก่อนเริ่มส่วนหลัง
\end{important}

\section{การตรวจสถานะและบันทึกการทำงาน}
\label{sec:monitor}

ตรวจสถานะของระบบด้วยคำสั่งในตารางที่~\ref{tab:monitor}
ส่วนหลังบันทึกคำขอทุกรายการและข้อผิดพลาดไว้ในบันทึกของคอนเทนเนอร์ \cmd{api}

\begin{table}[H]
\caption{การตรวจสถานะของระบบ}
\label{tab:monitor}
\centering
\begin{tabular}{|L{3.4cm}|L{\dimexpr\linewidth-7.6cm-6\tabcolsep-4\arrayrulewidth\relax}|L{4.2cm}|}
\hline
\textbf{สิ่งที่ตรวจ} & \textbf{คำสั่ง} & \textbf{ผลที่ปกติ} \\ \hline
สถานะบริการ & \cmd{docker-compose ps} & \cmd{db} และ \cmd{api} มีสถานะ \cmd{Up} \\ \hline
การตอบสนองของส่วนหลัง & \cmd{curl http://localhost:3001/api/health} & \cmd{\{"status":"ok"} \\ \hline
บันทึกการทำงาน & \cmd{docker-compose logs --tail=100 api} & ไม่มีข้อความ \cmd{Error} ซ้ำ ๆ \\ \hline
พื้นที่ดิสก์ & \cmd{df -h} และ \cmd{docker system df} & พื้นที่ว่างมากกว่าร้อยละ 20 \\ \hline
\end{tabular}
\end{table}

\section{ข้อควรระวังด้านความปลอดภัย}
\label{sec:security}

\begin{enumerate}
  \item ตั้งค่า \cmd{JWT_SECRET} และ \cmd{DB_PASSWORD} ใหม่ ห้ามใช้ค่าตัวอย่าง
  \item ตั้ง \cmd{NODE_ENV=production} ระบบจะไม่แสดงรายละเอียดข้อผิดพลาดภายในให้ผู้ใช้เห็น
  \item ลบหรือเปลี่ยนรหัสผ่านของบัญชีทดสอบ (ตารางที่~\ref{tab:demo-accounts})
  \item เปิด HTTPS ที่ Nginx เช่น ใช้ใบรับรองจาก Let's Encrypt
  \item ปิดพอร์ต 3001 และ 5432 จากภายนอกด้วยไฟร์วอลล์ ให้เข้าระบบผ่าน Nginx เท่านั้น
  \item คอนเทนเนอร์ที่ใช้ Compile ปิด shell escape และไม่มีการเชื่อมต่อเครือข่าย ห้ามเปลี่ยนค่าเหล่านี้
\end{enumerate}


% คู่มือสำหรับนักพัฒนา
\chapter{คู่มือสำหรับนักพัฒนา}
% =====================================================================
% บทที่ 7 คู่มือสำหรับนักพัฒนา (ผู้อ่าน: นักพัฒนาที่รับช่วงต่อ)
% =====================================================================
\label{ch:dev}

\begin{guide}
บทนี้สำหรับนักพัฒนาที่ต้องแก้ไขหรือพัฒนาระบบต่อ เช่น รุ่นน้องหรือหน่วยงานเจ้าของระบบ
เป้าหมายคือให้ผู้อ่านเริ่มพัฒนาได้เองโดยไม่ต้องถามผู้จัดทำ
\begin{itemize}
  \item สถาปัตยกรรมของระบบ และเทคโนโลยีที่ใช้พร้อมรุ่น
  \item โครงสร้างโฟลเดอร์ของซอร์สโค้ด และหน้าที่ของแต่ละส่วน
  \item การเตรียมเครื่องสำหรับพัฒนา การรันระบบ และการทดสอบ
  \item ฐานข้อมูล: ตารางหลักและความสัมพันธ์ (รายละเอียดทุกฟิลด์ใส่ภาคผนวก)
  \item API หรือส่วนต่อประสาน: สรุปหลักการในบทนี้ รายการเต็มใส่ภาคผนวก
  \item แนวทางพัฒนาต่อ: ขั้นตอนเพิ่มฟังก์ชันใหม่ ข้อตกลงในการเขียนโค้ด และการใช้ Git
  \item สิ่งที่ต้องทำก่อนนำรุ่นใหม่ขึ้นใช้งานจริง
\end{itemize}
แผนภาพ UML ฉบับเต็มอยู่ในบทที่ 4 ของเล่มรายงานแล้ว ให้อ้างถึงแทนการคัดลอกมาทั้งหมด \\
แอปบนมือถือ: เพิ่มการ build การ sign และการตั้งค่า endpoint ของส่วนหลัง \\
ระบบฮาร์ดแวร์/IoT: เพิ่มการกำหนดขา (pin) โปรโตคอลสื่อสาร และวิธีอัปโหลดเฟิร์มแวร์ \\
งานวิจัย/โมเดล AI: เพิ่มการเตรียมข้อมูล การฝึกโมเดลใหม่ ค่าพารามิเตอร์ และวิธีวัดผล
\end{guide}

\section{สถาปัตยกรรมของระบบ}
\label{sec:architecture}

ระบบแบ่งเป็น 3 ส่วนดังภาพที่~\ref{fig:architecture} ส่วนหน้าติดต่อส่วนหลังผ่าน REST API
และใช้ WebSocket สำหรับการแก้ไขเอกสารร่วมกัน ส่วนหลังสั่ง Compile เอกสารในคอนเทนเนอร์ TeX Live
ที่แยกจากระบบหลัก เทคโนโลยีที่ใช้แสดงดังตารางที่~\ref{tab:tech}

\begin{figure}[H]
\centering
\begin{tikzpicture}[
    node distance=18bp and 58bp,
    tier/.style={draw=cskkuStep, fill=cskkuStep!8, rounded corners=3pt, align=center,
      font=\normalfont\small, minimum height=50bp, text width=80bp, inner sep=3bp},
    aux/.style={tier, draw=black!55, fill=black!4},
    arr/.style={{Stealth[length=5bp]}-{Stealth[length=5bp]}, line width=0.8pt, draw=black!60},
    one/.style={-{Stealth[length=5bp]}, line width=0.8pt, draw=black!60},
    lbl/.style={font=\normalfont\footnotesize, text=black!75, inner sep=2bp, align=center}]
  \node[tier] (fe) {ส่วนหน้า\\React + Vite\\Monaco Editor};
  \node[tier, right=of fe] (be) {ส่วนหลัง\\Node.js + Express\\Yjs};
  \node[tier, right=of be] (db) {ฐานข้อมูล\\PostgreSQL\\(พัฒนา: SQLite)};
  \node[aux, below=of be] (tex) {Docker\\TeX Live + XeLaTeX};
  \draw[arr] ([yshift=8bp]fe.east) -- node[lbl, above]{REST /api} ([yshift=8bp]be.west);
  \draw[arr] ([yshift=-8bp]fe.east) -- node[lbl, below]{WebSocket\\/collab} ([yshift=-8bp]be.west);
  \draw[arr] (be) -- node[lbl, above]{Sequelize} (db);
  \draw[one] (be) -- node[lbl, right]{Compile} (tex);
\end{tikzpicture}
\caption{สถาปัตยกรรมของระบบ}
\label{fig:architecture}
\end{figure}

\begin{table}[H]
\caption{เทคโนโลยีที่ใช้ในระบบ}
\label{tab:tech}
\centering
\begin{tabular}{|L{3.6cm}|L{\dimexpr\linewidth-7cm-6\tabcolsep-4\arrayrulewidth\relax}|L{3.4cm}|}
\hline
\textbf{ส่วน} & \textbf{เทคโนโลยี} & \textbf{รุ่น} \\ \hline
ส่วนหน้า & React, Vite, Monaco Editor, Axios & 18 / 5 \\ \hline
การแก้ไขร่วมกัน & Yjs (CRDT), y-websocket, y-monaco & 13.6 \\ \hline
ส่วนหลัง & Node.js, Express, Sequelize & 20 / 4.18 / 6 \\ \hline
ฐานข้อมูล & PostgreSQL (ใช้งานจริง), SQLite (พัฒนา) & 15 / 3 \\ \hline
การยืนยันตัวตน & JWT และ bcrypt & -- \\ \hline
Compile เอกสาร & Docker, TeX Live, XeLaTeX & TeX Live 2023 \\ \hline
\end{tabular}
\end{table}

รายละเอียดการใช้งาน Node.js และ PostgreSQL ศึกษาเพิ่มเติมได้จากเอกสารของผู้พัฒนา~\cite{nodejs2025,postgres2025}

\section{โครงสร้างซอร์สโค้ด}
\label{sec:source-tree}

ซอร์สโค้ดแบ่งเป็นโฟลเดอร์ \filename{client} (ส่วนหน้า) และ \filename{server} (ส่วนหลัง)
โฟลเดอร์ที่ต้องแก้ไขบ่อยมีดังนี้

\begin{lstlisting}[style=cskkuterminal]
cskkuproject/
|-- client/src/
|   |-- App.jsx            เส้นทางของหน้าจอ (routing)
|   |-- pages/             หน้าจอหลัก
|   |-- components/        ส่วนประกอบที่ใช้ซ้ำ
|   |-- contexts/          สถานะการเข้าสู่ระบบ
|   `-- services/api.js    ฟังก์ชันเรียก API
|-- server/
|   |-- app.js             จุดเริ่มต้นของส่วนหลัง
|   |-- collaboration.js   WebSocket แก้ไขร่วมกัน
|   |-- config/            การเชื่อมต่อฐานข้อมูล
|   |-- models/            โมเดลข้อมูล (Sequelize)
|   |-- controllers/       การทำงานของแต่ละ API
|   |-- routes/            เส้นทาง API
|   |-- middleware/        ตรวจ token สิทธิ์ ข้อมูลนำเข้า
|   |-- templates/         แม่แบบ LaTeX ของโครงงานใหม่
|   |-- scripts/           สคริปต์ดูแลและทดสอบ
|   `-- docker/            Dockerfile ของ TeX Live
`-- docker-compose.yml
\end{lstlisting}

\section{การเตรียมเครื่องสำหรับพัฒนา}
\label{sec:dev-setup}

ระหว่างพัฒนา ส่วนหลังใช้ฐานข้อมูล SQLite (ไฟล์ \filename{server/database.sqlite}) เมื่อไม่ได้กำหนด
\cmd{DATABASE_URL} จึงไม่ต้องติดตั้ง PostgreSQL ซอฟต์แวร์ที่ต้องมีแสดงดังตารางที่~\ref{tab:dev-prereq}

\begin{table}[H]
\caption{ซอฟต์แวร์สำหรับเครื่องของนักพัฒนา}
\label{tab:dev-prereq}
\centering
\begin{tabular}{|L{3.8cm}|L{2.4cm}|C{2.6cm}|L{\dimexpr\linewidth-8.8cm-8\tabcolsep-5\arrayrulewidth\relax}|}
\hline
\textbf{ซอฟต์แวร์} & \textbf{รุ่น} & \textbf{จำเป็น} & \textbf{คำสั่งตรวจสอบ} \\ \hline
Node.js            & 18 ขึ้นไป     & \cmark & \cmd{node --version} \\ \hline
npm                & 9 ขึ้นไป      & \cmark & \cmd{npm --version} \\ \hline
Git                & 2.30 ขึ้นไป   & \cmark & \cmd{git --version} \\ \hline
TeX Live (XeLaTeX) & 2023 ขึ้นไป   & \cmark & \cmd{xelatex --version} \\ \hline
Docker             & 24 ขึ้นไป     & --     & \cmd{docker --version} \\ \hline
PostgreSQL         & 15 ขึ้นไป     & --     & \cmd{psql --version} \\ \hline
\end{tabular}
\end{table}

\begin{steps}
  \item ดาวน์โหลดซอร์สโค้ดและติดตั้งแพ็กเกจของทั้งสองส่วน
\begin{terminal}
git clone <repo-url> cskkuproject
cd cskkuproject/server && npm install
cd ../client && npm install
\end{terminal}
  \item สร้างไฟล์ตั้งค่าจากไฟล์ตัวอย่าง แล้วลบบรรทัด \cmd{DATABASE_URL} ออกเพื่อใช้ SQLite
\begin{terminal}
cd ../server
cp .env.example .env
\end{terminal}
  \item สร้างข้อมูลทดสอบ (ครั้งแรกเท่านั้น) และเริ่มส่วนหลังในเทอร์มินัลที่ 1
\begin{terminal}
npm run db:seed
npm run dev
\end{terminal}
  \item เริ่มส่วนหน้าในเทอร์มินัลที่ 2
\begin{terminal}
cd client
npm run dev
\end{terminal}
        \result{เปิด \url{http://localhost:5173} แล้วเข้าสู่ระบบด้วยบัญชีในตารางที่~\ref{tab:demo-accounts}}
\end{steps}

\begin{note}
ส่วนหน้าที่รันด้วย \cmd{npm run dev} ส่งต่อคำขอ \cmd{/api} และ \cmd{/collab} ไปยังพอร์ต 3001 ให้เอง
(ตั้งค่าไว้ใน \filename{client/vite.config.js}) ระหว่างพัฒนาจึงไม่ต้องตั้งค่า Nginx
\end{note}

\section{การทดสอบ}
\label{sec:dev-test}

สคริปต์ \filename{server/scripts/flowTest.js} ทดสอบการทำงานหลักผ่าน API จริง
ตั้งแต่สมัครสมาชิก สร้างและอนุมัติข้อเสนอ จนถึงการสร้างไฟล์และความเห็น
เริ่มส่วนหลังก่อน แล้วรันในเทอร์มินัลใหม่

\begin{terminal}
cd server
node scripts/flowTest.js
\end{terminal}

ผลลัพธ์แสดงเครื่องหมาย \cmark{} หน้ากรณีที่ผ่าน และสรุปจำนวนกรณีที่ผ่านและไม่ผ่านท้ายรายงาน

\begin{warning}
สคริปต์ทดสอบสร้างผู้ใช้และข้อมูลทดสอบลงฐานข้อมูล ใช้กับฐานข้อมูลสำหรับพัฒนาเท่านั้น
ล้างฐานข้อมูลสำหรับพัฒนาได้ด้วย \cmd{node scripts/resetDb.js} ซึ่งลบข้อมูลทั้งหมด
\end{warning}

\section{ฐานข้อมูล}
\label{sec:database}

ฐานข้อมูลมี 9 ตารางตามตารางที่~\ref{tab:db-tables} โมเดลของแต่ละตารางอยู่ใน \filename{server/models/}
และความสัมพันธ์กำหนดไว้ใน \filename{server/models/index.js}
ส่วนหลังสร้างตารางให้อัตโนมัติเมื่อเริ่มทำงาน และตั้งชื่อคอลัมน์แบบ snake\_case เช่น \cmd{created_at}

\begin{table}[H]
\caption{ตารางในฐานข้อมูล}
\label{tab:db-tables}
\centering
\begin{tabular}{|L{3.4cm}|L{\dimexpr\linewidth-8.4cm-6\tabcolsep-4\arrayrulewidth\relax}|L{5cm}|}
\hline
\textbf{ตาราง} & \textbf{ข้อมูลที่เก็บ} & \textbf{เชื่อมกับตาราง} \\ \hline
\cmd{users} & ผู้ใช้และบทบาท & -- \\ \hline
\cmd{proposals} & ข้อเสนอโครงงานและสถานะ & \cmd{users} (นักศึกษา อาจารย์) \\ \hline
\cmd{proposal_logs} & ประวัติการเปลี่ยนสถานะข้อเสนอ & \cmd{proposals}, \cmd{users} \\ \hline
\cmd{projects} & โครงงานที่อนุมัติแล้ว & \cmd{proposals} (หนึ่งต่อหนึ่ง) \\ \hline
\cmd{files} & ไฟล์ LaTeX ของโครงงาน เก็บเนื้อหาในฐานข้อมูล & \cmd{projects} \\ \hline
\cmd{comments} & ความเห็นในไฟล์ และการตอบกลับ & \cmd{files}, \cmd{users}, \cmd{comments} \\ \hline
\cmd{submissions} & การส่งงาน (สำเนาไฟล์ ณ เวลาส่ง) & \cmd{projects}, \cmd{users} \\ \hline
\cmd{deadlines} & กำหนดส่งงาน & \cmd{users} (ผู้สร้าง) \\ \hline
\cmd{notifications} & การแจ้งเตือนของผู้ใช้ & \cmd{users} \\ \hline
\end{tabular}
\end{table}

\section{API}
\label{sec:api}

ส่วนหลังให้บริการ REST API ที่ \cmd{/api} รับและส่งข้อมูลแบบ JSON รายการทั้งหมดอยู่ในภาคผนวก~ข
หลักการที่ใช้ร่วมกันทุก API มีดังนี้

\begin{enumerate}
  \item เข้าสู่ระบบด้วย \cmd{POST /api/auth/login} จะได้ token กลับมา
        แล้วส่ง token ในส่วนหัว \cmd{Authorization: Bearer <token>} ทุกคำขอ
  \item สิทธิ์ตามบทบาทตรวจด้วย middleware \cmd{roleGuard} ถ้าไม่มีสิทธิ์จะได้รหัส 403
  \item การแก้ไขร่วมกันเชื่อมต่อ WebSocket ที่ \cmd{/collab?room=project:<id>:file:<id>\&token=<token>}
  \item เมื่อเกิดข้อผิดพลาด ระบบตอบกลับด้วยรหัสสถานะ HTTP และข้อความในรูปแบบดังนี้
\end{enumerate}

\begin{terminal}
{ "error": { "code": 403, "message": "..." } }
\end{terminal}

ตัวอย่างการเรียก API ด้วย curl

\begin{terminal}
curl -X POST http://localhost:3001/api/auth/login \
  -H "Content-Type: application/json" \
  -d '{"email": "somchai.stud@kku.ac.th",
       "password": "password123"}'
\end{terminal}

\section{แนวทางการพัฒนาต่อ}
\label{sec:extend}

\subsection{การเพิ่มฟังก์ชันใหม่}

\begin{steps}
  \item เพิ่มหรือแก้โมเดลใน \filename{server/models/} และกำหนดความสัมพันธ์ใน \filename{models/index.js}
  \item เขียนการทำงานใน \filename{server/controllers/}
  \item เพิ่มเส้นทางใน \filename{server/routes/} พร้อม middleware \cmd{auth} และ \cmd{roleGuard}
        ตามสิทธิ์ (ตัวอย่างในโปรแกรมที่~\ref{lst:route}) แล้วลงทะเบียนเส้นทางใน \filename{app.js}
  \item เพิ่มฟังก์ชันเรียก API ใน \filename{client/src/services/api.js}
  \item เพิ่มหน้าหรือส่วนประกอบใน \filename{client/src/pages/} หรือ \filename{components/}
        และเพิ่มเส้นทางใน \filename{App.jsx}
  \item เพิ่มกรณีทดสอบใน \filename{scripts/flowTest.js} แล้วรันทดสอบตามหัวข้อ~\ref{sec:dev-test}
  \item ปรับคู่มือผู้ใช้และภาคผนวก~ข ให้ตรงกับฟังก์ชันใหม่
\end{steps}

\begin{lstlisting}[language=Java, caption={เส้นทาง API ที่จำกัดสิทธิ์ตามบทบาท (server/routes/deadlines.js)}, label={lst:route}]
const router = require('express').Router();
const ctrl = require('../controllers/deadlineController');
const auth = require('../middleware/auth');
const roleGuard = require('../middleware/roleGuard');

// ทุกบทบาทดูกำหนดส่งได้ แต่สร้างได้เฉพาะอาจารย์
router.get('/', auth, ctrl.list);
router.post('/', auth, roleGuard('instructor', 'advisor'), ctrl.create);

module.exports = router;
\end{lstlisting}

\subsection{ข้อตกลงในการพัฒนา}

\begin{enumerate}
  \item ชื่อตัวแปร ฟังก์ชัน และความเห็นในโค้ดเป็นภาษาอังกฤษ ข้อความที่แสดงต่อผู้ใช้เป็นภาษาไทย
  \item API ที่ต้องเข้าสู่ระบบใช้ middleware \cmd{auth} ทุกเส้นทาง และตรวจข้อมูลนำเข้าด้วยกฎใน
        \filename{middleware/validate.js}
  \item แยก branch ต่อหนึ่งฟังก์ชัน เช่น \cmd{feature/deadline-reminder} แล้วรวมเข้า \cmd{main}
        ผ่าน pull request หลังทดสอบผ่าน
  \item ห้าม commit ไฟล์ \filename{.env} และ \filename{database.sqlite}
\end{enumerate}

\section{การนำรุ่นใหม่ขึ้นใช้งานจริง}
\label{sec:deploy}

ขั้นตอนติดตั้งบนเครื่องแม่ข่ายอยู่ในหัวข้อ~\ref{sec:install-steps} และการปรับรุ่นอยู่ในหัวข้อ~\ref{sec:upgrade}
ก่อนส่งมอบรุ่นใหม่ให้ผู้ดูแลระบบ นักพัฒนาต้องตรวจรายการต่อไปนี้

\begin{enumerate}
  \item Build ส่วนหน้าด้วย \cmd{npm run build} ได้โดยไม่มีข้อผิดพลาด
  \item สคริปต์ทดสอบผ่านทุกกรณี
  \item ทดลองแก้ไขเอกสารร่วมกันจากเบราว์เซอร์ 2 เครื่องผ่าน Nginx เพื่อยืนยันว่า WebSocket ทำงาน
  \item เขียนบันทึกการเปลี่ยนแปลงของรุ่น โดยเฉพาะการเปลี่ยนโครงสร้างตาราง
  \item ปรับคู่มือเล่มนี้ และเพิ่มรายการในประวัติการแก้ไขเอกสาร
\end{enumerate}


% ผู้อ่านทุกกลุ่ม
\chapter{การแก้ไขปัญหาเบื้องต้น}
% =====================================================================
% บทที่ 8 การแก้ไขปัญหาเบื้องต้น (ผู้อ่าน: ทุกกลุ่ม)
% รวบรวมจากปัญหาที่พบจริงระหว่างทดสอบระบบกับผู้ใช้ (UAT) และระหว่างติดตั้ง
% =====================================================================
\label{ch:trouble}

\begin{guide}
รวบรวมจากปัญหาที่พบจริงระหว่างทดสอบระบบกับผู้ใช้ (UAT) และระหว่างติดตั้งบนเครื่องใหม่
\begin{itemize}
  \item ตารางอาการ--สาเหตุ--วิธีแก้ แยกปัญหาของผู้ใช้ออกจากปัญหาการติดตั้งและดูแลระบบ
  \item คัดลอกข้อความแจ้งเตือนให้ตรงกับที่ระบบแสดงจริง ผู้อ่านจะค้นหาข้อความนั้นในคู่มือได้
  \item คำถามที่พบบ่อย และช่องทางติดต่อขอความช่วยเหลือ พร้อมข้อมูลที่ควรแจ้งมา
\end{itemize}
\end{guide}

\section{ปัญหาที่ผู้ใช้พบบ่อย}

\begin{troubleshooting}[ปัญหาที่ผู้ใช้พบบ่อยและวิธีแก้ไข][tab:troubleshoot]
\problem{สมัครสมาชิกไม่ได้ ระบบแจ้งว่าอีเมลไม่ถูกต้อง}
        {ใช้อีเมลที่ไม่ได้ลงท้ายด้วย \cmd{@kku.ac.th}}
        {สมัครด้วยอีเมลของมหาวิทยาลัย}
\problem{Compile แล้วไม่ได้ไฟล์ PDF}
        {โค้ด \LaTeX{} ผิดรูปแบบ เช่น ปิด \cmd{\begin{...}} ไม่ครบ หรือวงเล็บปีกกาไม่ครบคู่}
        {อ่านบรรทัดที่ผิดใน COMPILE LOG แล้วแก้ไข}
\problem{ขึ้นข้อความ ``กำลังแปลผลอยู่''}
        {สมาชิกคนอื่นในโครงงานกำลัง Compile}
        {รอสักครู่แล้ว Compile ใหม่}
\problem{จุดสถานะที่แถบเครื่องมือเป็นสีแดง}
        {WebSocket ไม่ได้เชื่อมต่อ ระบบทำงานแบบออฟไลน์}
        {รีเฟรชหน้าเว็บ และบันทึกงานด้วย \keys{Ctrl+S}}
\problem{เห็นตัวอักษรเพิ่มขึ้นใน Editor เอง}
        {สมาชิกคนอื่นกำลังแก้ไขไฟล์เดียวกัน}
        {ดูรายชื่อผู้ที่ออนไลน์ที่แถบเครื่องมือ ระบบรวมการแก้ไขให้อัตโนมัติ}
\problem{ภาพที่แทรกไม่แสดงใน PDF}
        {ยังไม่ได้อัปโหลดไฟล์ภาพ หรือชื่อไฟล์ใน \cmd{\includegraphics} ไม่ตรง}
        {อัปโหลดภาพที่แฟ้มข้อมูล แล้วแก้ชื่อไฟล์ให้ตรงกัน}
\end{troubleshooting}

\section{ปัญหาการติดตั้งและดูแลระบบ}

\begin{troubleshooting}[ปัญหาการติดตั้งและดูแลระบบ][tab:troubleshoot-admin]
\problem{\cmd{EADDRINUSE: port 3001}}
        {มีโปรแกรมอื่นใช้พอร์ต 3001 อยู่}
        {ปิดโปรแกรมเดิม หรือเปลี่ยน \cmd{PORT} ในไฟล์ \filename{server/.env}}
\problem{\cmd{xelatex: command not found} หรือ Compile ไม่ได้ PDF ทุกโครงงาน}
        {ยังไม่ได้สร้างอิมเมจ TeX Live หรือไม่ได้ติดตั้ง TeX Live}
        {สร้างอิมเมจตามขั้นตอนที่ 3 ในหัวข้อ~\ref{sec:install-steps}}
\problem{\cmd{npm install} ล้มเหลวที่แพ็กเกจ SQLite}
        {เครื่องไม่มีเครื่องมือ build ของ native module}
        {Linux: \cmd{sudo apt install build-essential python3} \newline Windows: ติดตั้ง Visual Studio Build Tools}
\problem{เบราว์เซอร์แจ้ง \cmd{CORS error}}
        {\cmd{CLIENT_URL} ไม่ตรงกับที่อยู่ของส่วนหน้า}
        {แก้ \cmd{CLIENT_URL} ใน \filename{server/.env} แล้วเริ่มส่วนหลังใหม่}
\problem{ผู้ใช้แก้ไขพร้อมกันไม่ได้หลังติดตั้งบนเครื่องแม่ข่าย}
        {Nginx ไม่ได้ส่งต่อ WebSocket ที่ \cmd{/collab}}
        {เพิ่ม \cmd{proxy_set_header Upgrade} และ \cmd{Connection} ตามไฟล์ในหัวข้อ~\ref{sec:install-steps}}
\problem{\cmd{database is locked} (ระหว่างพัฒนา)}
        {มีหลายโปรแกรมเปิดไฟล์ SQLite พร้อมกัน}
        {ปิดส่วนหลังที่เปิดค้างทั้งหมด แล้วเริ่มใหม่}
\end{troubleshooting}

\section{คำถามที่พบบ่อย}

\begin{faq}
\question{ทำไมข้อเสนอโครงงานถูกส่งกลับ}
อาจารย์ที่ปรึกษาส่งกลับเพื่อให้แก้ไข ให้อ่านความเห็นของอาจารย์ที่หน้าข้อเสนอ แก้ไขตามความเห็น แล้วส่งใหม่

\question{เพิ่มรูปภาพในเอกสารอย่างไร}
อัปโหลดรูปที่แฟ้มข้อมูล แล้วใช้ปุ่ม \menu{แทรก > ภาพ} เพื่อแทรกคำสั่ง \cmd{\includegraphics}
จากนั้นแก้ชื่อไฟล์ในคำสั่งให้ตรงกับไฟล์ที่อัปโหลด

\question{อ้างอิงเอกสารอย่างไร}
เพิ่มรายการในไฟล์ \filename{references.bib} แล้วใช้คำสั่ง \cmd{\cite{key}} ในเนื้อหา

\question{ลืมรหัสผ่านต้องทำอย่างไร}
ติดต่อผู้ดูแลระบบตามช่องทางในหัวข้อ~\ref{sec:contact} ผู้ดูแลระบบจะตั้งรหัสผ่านชั่วคราวให้
(หัวข้อ~\ref{sec:users}) แล้วเปลี่ยนรหัสผ่านใหม่ที่หน้าโปรไฟล์หลังเข้าสู่ระบบ
\end{faq}

\section{การติดต่อขอความช่วยเหลือ}
\label{sec:contact}

หากพบปัญหาที่ไม่มีในคู่มือ ติดต่อผู้ดูแลระบบหรือผู้จัดทำได้ที่
สาขาวิชาวิทยาการคอมพิวเตอร์ วิทยาลัยการคอมพิวเตอร์ มหาวิทยาลัยขอนแก่น
โดยแจ้งชื่อผู้ใช้ เวลาที่พบปัญหา และข้อความแจ้งเตือนหรือภาพหน้าจอประกอบ


% ---------------------------------------------------------------------
% 4) ส่วนท้าย
% ---------------------------------------------------------------------
\printreferences     % ลบบรรทัดนี้ได้ถ้าคู่มือไม่มีการอ้างอิง

\appendix            % หน้าคั่น "ภาคผนวก" แล้วเริ่มภาคผนวก ก
\chapter{อภิธานศัพท์}
% =====================================================================
% ภาคผนวก ก อภิธานศัพท์ — คำเฉพาะที่ผู้อ่านอาจไม่คุ้นเคย เรียงตามตัวอักษร (อังกฤษก่อนไทย)
% =====================================================================
คำศัพท์เฉพาะที่ใช้ในคู่มือและในหน้าจอของระบบมีความหมายดังนี้

\begin{glossarylist}
\term{Compile}{การแปลงไฟล์ \LaTeX{} เป็นไฟล์ PDF}
\term{COMPILE LOG}{บันทึกผลการ Compile แสดงข้อความแจ้งข้อผิดพลาดและบรรทัดที่ผิด}
\term{Docker}{ซอฟต์แวร์สำหรับรันโปรแกรมในคอนเทนเนอร์ที่แยกจากระบบหลัก ใช้รันฐานข้อมูล ส่วนหลัง และ TeX Live}
\term{JWT (token)}{ข้อมูลยืนยันตัวตนที่ระบบออกให้หลังเข้าสู่ระบบ แนบไปกับทุกคำขอ}
\term{LaTeX}{ภาษาสำหรับจัดรูปแบบเอกสาร ใช้เขียนเล่มรายงานโครงงาน}
\term{Nginx}{เว็บเซิร์ฟเวอร์ที่แสดงส่วนหน้าและส่งต่อคำขอไปยังส่วนหลัง}
\term{Snapshot}{สำเนาไฟล์ทั้งหมดของโครงงาน ณ เวลาที่ส่งงาน}
\term{WebSocket}{ช่องทางสื่อสารแบบเรียลไทม์ระหว่างเบราว์เซอร์กับเครื่องแม่ข่าย ใช้ในการแก้ไขเอกสารร่วมกัน}
\term{ข้อเสนอโครงงาน}{เอกสารขออนุมัติหัวข้อโครงงาน ประกอบด้วยชื่อโครงงาน สมาชิก และอาจารย์ที่ปรึกษา}
\term{ฉบับร่าง}{ข้อเสนอที่บันทึกไว้แต่ยังไม่ได้ส่งให้อาจารย์ที่ปรึกษา แก้ไขได้ตลอด}
\term{หัวหน้าโครงงาน}{นักศึกษาผู้สร้างข้อเสนอโครงงาน}
\end{glossarylist}


\chapter{รายการ API ของระบบ}
% =====================================================================
% ภาคผนวก ข รายการ API — ส่วนอ้างอิงของคู่มือสำหรับนักพัฒนา (บทที่ 7)
% ใช้ longtable เมื่อรายการยาวเกิน 1 หน้า (ชื่อตารางพิมพ์ซ้ำพร้อม "(ต่อ)" ให้อัตโนมัติ)
% ทุกเส้นทางขึ้นต้นด้วย /api ยกเว้นที่ระบุไว้ สิทธิ์ "ผู้ใช้" = ผู้ที่เข้าสู่ระบบแล้วทุกบทบาท
% =====================================================================
\label{app:api}

\begin{guide}
ภาคผนวกเก็บข้อมูลอ้างอิงที่ยาวเกินจะใส่ในบท เช่น รายการ API ทั้งหมด พจนานุกรมข้อมูล (ทุกตาราง ทุกฟิลด์)
หรือรหัสข้อผิดพลาด จัดเป็นตารางเรียงตามกลุ่ม และปรับทุกครั้งที่ระบบเปลี่ยน
\end{guide}

ระบบรุ่น \systemversion{} มี REST API 37 รายการ แบ่งตามกลุ่มดังตารางที่~\ref{tab:api}
สิทธิ์ ``ผู้ใช้'' หมายถึงผู้ที่เข้าสู่ระบบแล้วทุกบทบาท

\begin{longtable}{|C{1.7cm}|L{5.4cm}|L{2.9cm}|L{\dimexpr\linewidth-10cm-8\tabcolsep-5\arrayrulewidth\relax}|}
\caption{รายการ API ของระบบ}\label{tab:api}\\
\hline
\textbf{Method} & \textbf{เส้นทาง} & \textbf{สิทธิ์} & \textbf{การทำงาน} \\ \hline
\endfirsthead
\caption[]{รายการ API ของระบบ~\cskkuContinued}\\
\hline
\textbf{Method} & \textbf{เส้นทาง} & \textbf{สิทธิ์} & \textbf{การทำงาน} \\ \hline
\endhead
\multicolumn{4}{|l|}{\textbf{การยืนยันตัวตน} (\cmd{/api/auth})} \\ \hline
POST & \cmd{/register} & ทุกคน & สมัครสมาชิก \\ \hline
POST & \cmd{/login} & ทุกคน & เข้าสู่ระบบ รับ token \\ \hline
GET & \cmd{/me} & ผู้ใช้ & ข้อมูลผู้ใช้ปัจจุบัน \\ \hline
PUT & \cmd{/me} & ผู้ใช้ & แก้ไขข้อมูลส่วนตัว \\ \hline
PUT & \cmd{/password} & ผู้ใช้ & เปลี่ยนรหัสผ่าน \\ \hline
GET & \cmd{/advisors} & ผู้ใช้ & รายชื่ออาจารย์ \\ \hline
GET & \cmd{/students} & ผู้ใช้ & รายชื่อนักศึกษา \\ \hline
\multicolumn{4}{|l|}{\textbf{ข้อเสนอโครงงาน} (\cmd{/api/proposals})} \\ \hline
GET & \cmd{/} & ผู้ใช้ & รายการข้อเสนอ \\ \hline
POST & \cmd{/} & นักศึกษา & สร้างข้อเสนอ \\ \hline
GET & \cmd{/:id} & ผู้ใช้ & รายละเอียดข้อเสนอ \\ \hline
PUT & \cmd{/:id} & นักศึกษา & แก้ไขข้อเสนอ (ฉบับร่างหรือส่งกลับ) \\ \hline
PUT & \cmd{/:id/submit} & นักศึกษา & ส่งข้อเสนอ \\ \hline
PUT & \cmd{/:id/review} & อาจารย์ที่ปรึกษา & อนุมัติหรือส่งกลับแก้ไข \\ \hline
GET & \cmd{/:id/logs} & ผู้ใช้ & ประวัติการเปลี่ยนสถานะ \\ \hline
\multicolumn{4}{|l|}{\textbf{โครงงานและไฟล์} (\cmd{/api/projects})} \\ \hline
POST & \cmd{/} & นักศึกษา & สร้างโครงงาน \\ \hline
GET & \cmd{/:id/files} & ผู้ใช้ & รายการไฟล์ \\ \hline
PUT & \cmd{/:id/files/:fid} & นักศึกษา อาจารย์ที่ปรึกษา & บันทึกไฟล์ \\ \hline
POST & \cmd{/:id/files/new} & นักศึกษา อาจารย์ที่ปรึกษา & สร้างไฟล์ใหม่ \\ \hline
POST & \cmd{/:id/files} & นักศึกษา อาจารย์ที่ปรึกษา & อัปโหลดไฟล์ \\ \hline
DELETE & \cmd{/:id/files/:fid} & นักศึกษา อาจารย์ที่ปรึกษา & ลบไฟล์ \\ \hline
POST & \cmd{/:id/compile} & นักศึกษา อาจารย์ที่ปรึกษา & Compile เป็น PDF \\ \hline
GET & \cmd{/:id/pdf} & ผู้ใช้ & ดาวน์โหลด PDF \\ \hline
\multicolumn{4}{|l|}{\textbf{ความเห็น} (\cmd{/api})} \\ \hline
GET & \cmd{/projects/:id/comments} & ผู้ใช้ & รายการความเห็นของโครงงาน \\ \hline
POST & \cmd{/projects/:id/comments} & ผู้ใช้ & เพิ่มความเห็น \\ \hline
POST & \cmd{/comments/:id/reply} & ผู้ใช้ & ตอบกลับความเห็น \\ \hline
PUT & \cmd{/comments/:id/resolve} & ผู้ใช้ & ปิดความเห็น (แก้ไขแล้ว) \\ \hline
\multicolumn{4}{|l|}{\textbf{การส่งงาน} (\cmd{/api/submissions})} \\ \hline
POST & \cmd{/} & ผู้ใช้ & ส่งงาน \\ \hline
GET & \cmd{/} & ผู้ใช้ & รายการการส่งงาน \\ \hline
PUT & \cmd{/:id} & ผู้ใช้ & ตรวจงาน \\ \hline
\multicolumn{4}{|l|}{\textbf{การแจ้งเตือน} (\cmd{/api/notifications})} \\ \hline
GET & \cmd{/} & ผู้ใช้ & รายการแจ้งเตือน \\ \hline
PUT & \cmd{/:id/read} & ผู้ใช้ & อ่านแล้ว \\ \hline
PUT & \cmd{/read-all} & ผู้ใช้ & อ่านแล้วทั้งหมด \\ \hline
\multicolumn{4}{|l|}{\textbf{แดชบอร์ดและกำหนดส่ง} (\cmd{/api/dashboard}, \cmd{/api/deadlines})} \\ \hline
GET & \cmd{/dashboard/summary} & ผู้ดูแลรายวิชา อาจารย์ที่ปรึกษา & สรุปจำนวนโครงงานและการส่งงาน \\ \hline
GET & \cmd{/dashboard/export} & ผู้ดูแลรายวิชา อาจารย์ที่ปรึกษา & ส่งออกไฟล์ CSV \\ \hline
GET & \cmd{/deadlines} & ผู้ใช้ & รายการกำหนดส่ง \\ \hline
POST & \cmd{/deadlines} & ผู้ดูแลรายวิชา อาจารย์ที่ปรึกษา & สร้างกำหนดส่ง \\ \hline
\multicolumn{4}{|l|}{\textbf{อื่น ๆ}} \\ \hline
GET & \cmd{/api/health} & ทุกคน & ตรวจสถานะส่วนหลัง \\ \hline
WS & \cmd{/collab?room=...&token=...} & ผู้ใช้ & แก้ไขไฟล์ร่วมกัน (Yjs) \\ \hline
\end{longtable}


\end{document}
